\documentclass[a4paper,11pt,openany]{ltjsbook}
\usepackage{luatexja}
\usepackage{luatexja-fontspec}
\usepackage{amsmath,amssymb,amsfonts,amsthm}
\usepackage{geometry}
\geometry{top=24mm,bottom=24mm,left=24mm,right=24mm,footskip=10mm}
\usepackage{graphicx}
\usepackage{tcolorbox}
\tcbuselibrary{skins,breakable}
\usepackage{hyperref}
\usepackage{enumitem}
\usepackage{booktabs}
\usepackage{fancyhdr}

\definecolor{primaryblue}{RGB}{20, 60, 110}
\definecolor{softblue}{RGB}{238, 244, 252}
\definecolor{darkslate}{RGB}{35, 45, 60}
\definecolor{warmgold}{RGB}{185, 125, 20}
\definecolor{softgold}{RGB}{254, 249, 235}
\definecolor{accentcrimson}{RGB}{160, 30, 45}

\hypersetup{
    colorlinks=true,
    linkcolor=primaryblue,
    citecolor=darkslate,
    urlcolor=accentcrimson,
    pdftitle={特殊および一般相対性理論について —— 原著完全直訳・数学校閲版},
    pdfauthor={アルベルト・アインシュタイン}
}

\pagestyle{fancy}
\fancyhf{}
\fancyhead[LE,RO]{\thepage}
\fancyhead[RE]{\leftmark}
\fancyhead[LO]{\rightmark}
\renewcommand{\headrulewidth}{0.4pt}

\title{
    {\huge\bfseries\gtfamily 特殊および一般相対性理論について}\\[12pt]
    {\Large\gtfamily —— 原典完全直訳・数学校閲・学術決定版 ——}\\[18pt]
    {\large 原題：\textit{Relativity: The Special and the General Theory} (1916/1920)}\\[6pt]
    {\normalsize Project Gutenberg eBook \#30155}
}
\author{\textbf{アルベルト・アインシュタイン} 著\\[4pt]教育情報分析研究会 訳・数学校閲}
\date{2026年10月 完全決定版}

\begin{document}

\maketitle

\frontmatter

\chapter*{訳者まえがき：アインシュタイン自筆の普及的名著}
\addcontentsline{toc}{chapter}{訳者まえがき：アインシュタイン自筆の普及的名著}

アルベルト・アインシュタイン（Albert Einstein, 1879--1955）が自ら一般の読者に向けて執筆した唯一の本格的解説書『特殊および一般相対性理論について（原題：\textit{Über die spezielle und die allgemeine Relativitätstheorie, gemeinverständlich}）』は、1916年に初版が出版され、1920年にロバート・W・ローソン（Robert W. Lawson）によって英語に翻訳されました（Project Gutenberg eBook \#30155）。

一般にアインシュタインの論文といえば難解なテンソル計算や重力方程式が想起されますが、本書におけるアインシュタインの筆致は驚くほど平易かつ誠実です。彼は序文においてルートヴィヒ・ボルツマンの格言「エレガンス（優雅さ）の追求は仕立屋と靴屋に任せるべきである」を引用し、技巧的なレトリックを排して、概念の必然性と物理的意味を徹底して明快に説き明かしています。

本全訳・数学校閲版では、一切の要約や省略を行わず、アインシュタインの原著英文（全32節および付録4篇）の論理構造と数理的展開を完全に日本語へ直訳いたしました。さらに、各節の数式と時空変換の数理的整合性を厳密に確認・校閲し、原典の論理を正確に追体験できる構成としております。

\clearpage
\tableofcontents
\clearpage

\mainmatter


\clearpage
\chapter*{原著序文（Preface）}
\addcontentsline{toc}{chapter}{原著序文（Preface）}

本書は、理論物理学の数学的装備に精通していないものの、科学および哲学的な観点から相対性理論に興味のある読者に向けて、可能な限り正確な洞察を提供することを目的としている。この作品は、大学入試レベルの教育水準を持つ読者に対して書かれているが、本は短いため、読者は相当な忍耐と意志力を発揮する必要がある。著者は、主なアイデアを最も単純で理解しやすい形で提示しようと努め、そして大抵の場合、実際の順序と連続性においてそれらが生まれた順に提示した。明瞭さのためには、私は自己を繰り返しの頻度から解放せざるを得なかったように思われ、そのような点は Presentation の elegance に関心を払う必要がないと、優れた理論物理学者 L. ボルツマンの教訓に従った。私は読者から本質的な難しさを隠すことを主張していないが、一方で、物理学的基礎を「半ば冷遇」したので、物理学に不慣れな読者は木から森を見ることができるように感じないだろう。本書は誰かに数時間の誘発的な思考を与えることを願っている！


\clearpage
\part{特殊相対性理論（The Special Theory of Relativity）}



\clearpage
\chapter{第I節：Physical Meaning of Geometrical Propositions}

I.
幾何学的命題の物理学的意味


あなたたちは学校時代に、この本を読むほとんどの人々がユークリッド幾何学の壮麗な建築物と出会い、その偉大なる構造について思い出すでしょう—無数の時間、慎重な教師たちによって追われて、高層の階段を駆け登るという経験は忘れられません。過去の経験から、これらの科学の最も偏った命題すら偽りであると述べる者は誰も蔑むべき存在であると考えるでしょう。しかしもし有人が問うた場合を考えましょう：「それでは、これらの命題が真実であるという主張は何か意味するのか？」この問いに対して少し考えてみましょう。

幾何学は「平面」、「点」、「直線」といった概念から出発します。これらには多少の明確なイメージを連想させますが、その上に単純な命題（公理）が置かれ、これらの概念に基づいて真実であると受け入れられることが好ましいと考えられます。それらの公理に基づき論理的な過程によって、残りのすべての命題は示されるべきであり、つまり証明されます。ある命題が正しく（真実である）とは、それが認知された方法から公理を導き出すことにより定義されます。個々の幾何学的命題の「真実性」は、公理の「真実性」に帰着されるのです。

しかし既に長い間知られているように、最後の問いは幾何学の手法では回答不可能であり、またその問い自体には意味がないことが明らかになっています。直線が二点を通るかどうかを問うことはできない。我々はただ、「直線」というものを扱っていることに注意し、それら全てがその上に位置する二点によって一意に定められるという性質を付与することだけが幾何学の対象であると言えるでしょう。「真実」の概念は純粋な幾何学の主張と一致しないのです。なぜなら「真実」という言葉は最終的には現実の物体との対応関係を示すものとして使われることが多いからです。しかし幾何学はその概念が経験的対象に対してどのような関係にあるかという問題ではなく、これらの概念間の論理的な接続のみに関心を持つのです。

これは幾何学の命題が「真実」であると感じられる理由を理解するのが難しくない。幾何学的な概念は自然の中のより或いはあまり正確な対象に応答しており、これらはこれらの概念が生まれる唯一の原因であることは明らかだ。幾何学はその構造に最大限の論理的統一性を与えるためにこのような道を避けるべきである。

例えば、「距離」という概念は実質的に剛体上の二つの明確な位置を見ることで形成されることが深く習慣化されている。さらに、私たちは三点が一直線上にあると見做すことに慣れており、観察のために片目を使ってその位置を一致させることができ、適切な視点を選べばそれが可能である。

この思考の習慣に従って、ユークリッド幾何学の命題に単一の命題「実質的に剛体上の二つの点は位置が変化しても常に同じ距離（線間隔）に対応する」を追加すれば、ユークリッド幾何学の命題は実質的に剛体の相対的な位置に関する命題に分解される。このような補完を受けた幾何学は物理学の一分野として扱われるべきである。

このように解釈された幾何学的命題の「真実性」について正当に問うことができる。なぜなら、これらの命題が我々が幾何学的概念に関連付けた現実的な対象で満足しているかどうかを問うことが正当化されるからだ。より正確な表現では、このような意味での幾何学的命題の「真実性」は規則とコンパスによる構築に対するその有効性であることを理解することができる。

なお、自然の対象も一直線に対応付けられる。剛体上の三点 $A, B$ および $C$ が一直線上にあるとは、点 $A$ と $C$ が与えられたとき、点 $B$ を選んでその距離 $AB$ と $BC$ の和が最小になる状態を指す。この不完全な提案は現在の目的には十分である。

当然、幾何学的命題の「真理性」はこの意味での信憑性は不完全な経験にのみ基礎を置いているものである。現時点では、「真理性」を仮定するが、後段（一般相対論の理論において）でこれは「真理性」が制限されていることを確認し、その制限の範囲について検討することにする。


\clearpage
\chapter{第II節：The System of Co-ordinates}

II.
座標系の体系


物理的な距離の解釈に基づいて、剛体上の二点間の距離を測定によっても定義できる位置にあります。この目的のために、「距離」（棒 $S$）が必要で、一度と決めてその標準として使用します。もし、現在、$A$ および $B$ が剛体上の二点であれば、幾何学的な規則に従ってそれらの点を結ぶ線を作ることができます。その後、$A$ から $S$ の距離を一定回数繰り返し測定して $B$ まで到達します。これらの操作の回数が距離 $AB$ の数値的な尺度となります。これが長さの測定の基礎です。

 ここで、測定結果が整数となることを仮定しています。つまり、測定が整数を残す場合の問題は、分割された測量棒を使用することで解決されます。その導入は基本的な新しい方法を必要としない。

事件や物体の場面の説明または空間上の位置の記述は、参照する剛体（参考系）上の点を特定することに基づいています。これは科学的な記述だけでなく、日常生活でも同様です。「トラファルガー広場、ロンドン」のような場所指定を分析すると、次の結果が得られます。地球が参照する剛体であり、「トラファルガー広場、ロンドン」は特定の点で、その名前が付けられており、空間上事件や物体と一致しています。

 私は「ポツダム広場、ベルリン」というドイツ語版に引用されているものよりも、「トラファルガー広場、ロンドン」の方が英語の読者にとって馴染みがあると考えたためこの選択をしました。（R. W. L.）

 「空間上の一致」の意味についてさらに調査する必要はありません。この概念は実際での適用性がほとんど疑問視されないほど明白であると解釈できます。

この原始的な場所指定法は、堅固な物体の表面に存在する場所のみを扱い、その表面に他のものと区別できる点が存在することに依存しています。しかし、これらの制限から解放されるためには、場所の指定自体の性質を変えずに両方とも克服することができます。たとえば、トラファルガー広場上空に浮かんでいる雲がある場合、地球の表面に対する雲の位置を特定するために、広場に垂直に立った柱を立てることができます。この柱が雲に達する長さを標準測量棒で計測し、柱の底の位置を指定することで、完全な場所指定が得られます。この例に基づいて、位置概念の精密化の方法を見ることができます。

(a) 我々は、場所指定が参照される堅固な物体に補完的な形で想像し、その補完された物体が求められる物体に達するようにします。

(b) 物体の位置を特定するために、指定点ではなく（ここでは測量棒で柱の長さ）数値を使用します。

(c) 雲の高さについて話す場合でも、雲に達する柱が立てられていなくても、地面からの異なる位置から雲に対する光学観測を行い、光の伝播特性を考慮に入れることにより、雲に達するために必要な柱の長さを求めることができます。

この考慮から、位置の説明において、名前を持つ特定の位置点（参照基準体上の）が存在することに依存せず、数値的な尺度を使用して自己独立化することが有利であることがわかります。測量物理学では、デカルト座標系の適用によってこの目標を達成します。

以下の文章を翻訳しました：

この装置は、互いに直角を成す三つの平面の表面で構成され、堅固な基準体に固定されています。座標系において、ある事件の場面を特定する（主として）には、その事件が存在する場所から三つの垂直線または座標（$x, y, z$）を下ろす長さを指定すればよいです。これらの三つの垂直線の長さは、ユークリッド幾何学に従って規則と方法を定めた操作を行うことにより測定することができます。

実際には、この座標系を構成する堅固な平面が一般的には利用できませんし、さらに座標の大きさは剛体の尺を使用した直接的な測定ではなく、間接的な手段で決定されます。物理学および天文学者の結果が明瞭性を保つためには、位置の指定の物理的意味は常に上記の考慮事項に従って探求する必要があります。

以上の観点において、以下の補足があります：これらの見解の精密化と修正は、本書の第二部で一般相対論に取り組む際に必要となります。

したがって我々は以下のように結論付けることができます：空間における事件の描写には常に、その事件を参照するために堅固な体を使用する必要があります。その結果得られる関係性では、ユークリッド幾何学の法則が「距離」について成り立つと仮定されています。「距離」とは、剛体上の二つのマークによる物理的な規定によって表現されます。


\clearpage
\chapter{第III節：Space and Time in Classical Mechanics}

III.

古典力学における空間と時間

力学の目的は、物体が「時間」とともに空間での位置をどのように変化するかを描写することである。このように述べるのは、明瞭性に対する神聖な精神に重罪を犯すようなものであり、真剣な反省と詳細な説明なしでその目的を述べることは慎むべきだ。では、これらの罪を明らかにしてみよう。

ここでの「位置」と「空間」が何を指しているのかは不明である。私は列車の窓から、一様に移動する列車の中で石を放り投げ、路肩に落ちるのを見ている。空気抵抗を無視すると、私は石が一直線に下落することを見る。路肩でその行為を目撃した歩行者は、石が円錐形の軌道を描いて地上に落下していることを観察する。では、このように石が通過する「位置」は「実際には」一直線上にあるのか、あるいは円錐形の軌道上にあるのか？さらに、「空間での運動」とは何を指すか？前節の考慮から答えは明らかである。まず第一に、我々は曖昧な言葉「空間」を使わず、実際には認識できないことを正直に認め、それを「参照系に対する相対的な移動」と置き換える。前節で詳細に定義した参照系（列車の窓や路肩）に対する位置を再確認することができる。代わりに「参照系」を「座標系」という数学的に描写するための有用な概念として挿入すると、次のように述べることができる：石は、列車に固定された座標系に対して一直線を描くが、地面（路肩）に固定された座標系に対しては円錐形の軌道を描く。この例を通じて明らかになるように、独立して存在する「軌跡」（軌道曲線）などというものはなく、ある参照系に対する相対的な軌跡のみが存在する。

以下の運動を完全に記述するためには、物体が位置を時間とともにどのように変化するかを指定する必要があります；すなわち、軌道上の各点について、その点に物体が存在しているのはどの時刻であるかを述べる必要があります。これらのデータは、この定義によってこれらの時刻の値が観測可能な測定結果として本質的に扱えるような時間の定義で補足されるべきです。古典力学の立場から見れば、我々は以下の方法でこの要件を満たすことができます。二つの同一構造の時計を想定します；鉄道車窓の男は一つを持ち、歩道の男はもう一つを持っています。各観測者は、自分の手に持っている時計の ticking ごとに石がどの位置にあるかを彼の参照体上での位置として決定します。ここで、光の伝播速度の有限性によって生じる不正確さについては考慮していません。また、ここに存在する二つ目の困難についても後で詳しく扱います。


\clearpage
\chapter{第IV節：The Galileian System of Co-ordinates}

IV. ガリレオ・ニュートンの力学系

ガリレオ＝ニュートン力学の基本法則、すなわち慣性の法則は、次のように述べることができる：十分に他の物体から離れた物体は、静止状態か一直線上の等速運動の状態を維持する。この法則は、物体の運動だけでなく、機械的な説明で使用できる許容される参照系や座標系も示唆している。視覚的に固定された星々は、慣性の法則が高精度で成り立つと考えられる物体である。もし地球に固着した剛体座標系を使用する場合、この系に対する相対的な観測では、各恒星は天文学的1日間にかけて巨大な円周を描くことがわかる。これは慣性の法則の主張と矛盾している結果である。従って、この法則に従うなら、これらの運動は固定星が円周上に動かない座標系に対してのみ説明できる。慣性の法則がその相対的な静止状態を有する座標系は「ガリレオ系」と呼ばれる。ガリレオ＝ニュートン力学の法則は、ガリレオ系においてのみ有効と見なされるべきである。


\clearpage
\chapter{第V節：The Principle of Relativity (in the Restricted Sense)}

第五章

相対性原理（狭義）

最も大きな明瞭さを得るため、私たちは列車が一定速度で移動しているという例に戻りましょう。その運動を「一定速度の平滑な移動」と呼びます（「一定速度」は常に同じ速度と方向であることを意味し、「平滑な移動」は列車が路肩に対して位置が変化するが回転していないことを意味します）。空を飛ぶ Raven が、路肩から観察したときの運動が一定速度で一直線であるように想像しましょう。もし列車の中から Raven の動きを観察すると、その動きは速度と方向が異なるものの、依然として一定速度で一直線上を移動していることがわかります。抽象的に表現すれば：座標系 $K$ に対して一定速度の直線運動をしている質量 $m$ について、座標系 $K′$ が $K$ に対する一定速度の平滑な移動であるとき、$K′$ に対してその質量も一定速度で一直線上を移動していることが成り立ちます。前節で述べたとおり、これは以下の結論を導きます：

もし $K$ がガリレオ座標系であるなら、関係してくる座標系 $K′$ も $K$ に対する一定速度の平滑な移動であるとき、ガリレオ座標系です。つまり、$K′$ に対してガリレーニュートンの力学法則は $K$ に対して同様に成り立ちます。

さらに一般化するためには、以下のように述べることができます：もし $K$ に対する $K′$ が一定速度で移動し、回転がない座標系であるなら、自然現象は $K′$ 上でも $K$ 上と同様の一般的な法則に従って進行します。この主張は狭義の相対性原理と呼ばれます。

古典力学の助けですべての自然現象を表現できるという信念がある限り、相対性原理の妥当性を疑う必要はなかった。しかし、最近の電磁気学と光学の発展を考えると、古典力学がすべての自然現象に対する物理学的な説明の基礎として不十分であることが明らかになった。この局面で相対性原理の妥当性について議論する時期が成熟し、その答えが否定的である可能性も考えられなくなった。

それでも、二つの一般的な事実は、相対性原理の妥当性を支持する上で非常に有力だ。古典力学が十分に広範でないとしても、天体の運動については驚くほど詳細な正確さで予測できるため、「真実」の一端を提供していることは認められるべきである。従って、相対性原理は機械学の領域において極めて正確に適用されるべきである。しかし、そのような広範囲に及ぶ原理が一つの現象領域では非常に正確に成り立つのに、別の領域では無効であるとは、事前に見当違いであると推測される。

我々は今や第二の論点へと進もう。これについては後に再び立ち返ることになる。もし相対性原理（制限された意味での）が成り立たないなら、互以匀速運動するガリレオ座標系 $K, K′, K″$, などは自然現象の説明において等価でなくなる。その場合、我々は自然法則が特に単純な形で表現可能であると信じざるを得ないし、もちろんそれは特定の一連の運動状態（$K0$）を参照基準体として選んだ場合に限られる。この系 $K0$ を「絶対的に静止している」と呼ぶことが正当化され、他のガリレオ座標系 $K$ は「動いている」と呼び得るだろう。

たとえば、我々の岸壁が $K0$ の場合、我々の鉄道車両は $K$ の参照基準体となり、この系に対する自然法則は $K0$ よりも単純ではない。その簡略化は、車両 $K$ が $K0$ 相対的に動いている（すなわち「実際には」）からである。一般の自然法則において $K$ を参照基準体として用いた場合の定式化では、車両の速度の大きさと向きは必然的に重要な役割を果たすことになる。例えば、方向に平行に配置されたオルガン管で発する音と、その軸が垂直に配置されたときの音とは異なることを期待せざるを得ないだろう。

現在、地球が太陽繞る軌道運動によって、地球は約30キロメートル/秒の速度で鉄道路車に比定される。相対性原理が成り立たなければ、我々はその瞬間における地球の運動方向が自然法則に含まれていることや、物理系の行動が地球に対する空間の向きによって依存することを期待するだろう。これは、一年間に地球の回転速度の方向が変化することから、地球は一年中相対的な静止状態で仮定された系統$K0$に対しては存在しないという観測結果から、地球は一年を通じて相対的に静止しているわけではないことを意味する。しかし、地球上の物理空間には、異なる方向に対する物理的非等価性のような不安定な特性が見出されることはない。これは相対性原理を支持する有力な証拠である。


\clearpage
\chapter{第VI節：The Theorem of the Addition of Velocities employed in Classical Mechanics}

VI.

古典力学において用いられる速度の加法定理

列車 carriage が鉄道 rails 上を一定の速度 $v$ で走行していると仮定し、その中で人間が旅行方向沿いに長さを覆う速度 $w$ で移動する場合、人間は1秒間にどれだけ進むか、言い換えれば、人間の相対速度 $W$ は何であるか？唯一の可能な答えは以下の考慮から得られる：もし人間が1秒間立ち止まっていたら、彼は列車の速度 $v$ に等しい距離 $v$ を鉄道に対して進むことになる。しかし、彼が歩行すると、列車に対してさらに距離 $w$ を覆うのであり、したがって鉄道に対しても同じ距離 $w$ を覆う。つまり、1秒間において人間は速度 $W = v + w$ で鉄道に対して進むことになる。後ほど見るように、この結果は古典力学において用いられる速度の加法定理を表し、これは維持できないことが明らかになる；言い換えれば、我々が書いた法則は現実には成り立たない。しかし現在はその正確性を仮定する。


\clearpage
\chapter{第VII節：The Apparent Incompatibility of the Law of Propagation of Light with the Principle of Relativity}

VII.

光の伝播に関する簡明な法則と相対性原理との表面上の不一致

物理学においては、空気中で光が伝播する法則はほとんど単純なものである。学校に通う子供たちはみな、あるいは少なくとも信じているが、この伝播は一直線に速度 $c$ = 300,000 km./sec. で行われると知っている。実際には、光の色によってこの速度が異なるわけではなく、固定星がその暗い隣接星体によって食された際に、各色ごとの発射の最小値が同時に観測されない場合があるから、我々は非常に正確にこの速度がすべての色に対して同じであることを知っている。このような考慮に基づいて観測された双星に関する類似の結論により、オランダの天文学者デ・サイターもまた、光の伝播速度が光源の運動速度によって依存しないことも示した。伝播速度が「空間」における方向に依存すると仮定するのは、自体において不確かなものである。

つまり、学校で子供たちは簡明な法則として光の速度 $c$（真空の場合）が一定であると信じているとする。誰が想像するだろうか、この単純な法則が真面目に考えを重ねる物理学者たちに最大の知的困難をもたらすとは。これらの困難はどのようにして生じるのかを考えてみよう。

当然、光の伝播過程（そして他のすべての過程も同様）は剛体の参照系（座標系）に参照しなければなりません。そのような系として列車線路上の軌道を再度選択します。空気を除去したと仮定すると、光線が軌道上を伝播するとき、軌道に対して光線の先端は速度 $c$ で伝播することを見ることができます。

今度は、列車が鉄道路線上を速度 $v$ で移動しているとします。その方向は光線と同じですが、当然その速度は非常に小さいです。列車上での光線の伝播速度を求めましょう。前節で考慮した考え方に基づくことは明らかであり、光線は列車上の歩行者に相当します。軌道に対する歩行者の速度 $W$ はここでは光線の速度が軌道に対して速度 $c$ で伝播することを表します。求める列車に対する光線の速度を $w$ とすると、

\[
w = c - v
\]

となります。

したがって、列車上での光線の伝播速度は $c$ よりも小さくなることが分かります。

しかし、この結果は第五節で述べた相対性原理と矛盾します。例えば、すべての自然法則は原則として相対性原理に従い、鉄道路線上を参照系とするときと同じように列車上の参照系でも光線の伝播法則が同じであるべきです。しかし、上記の考慮からこれは不可能なようです。もし全ての光線が軌道に対して速度 $c$ で伝播するとすれば、これにより必ずしも別の光線の伝播法則が列車上で成り立つことが必要であり、これが相対性原理に反する結果となります。

この難局においては、相対性原理を放棄するか、真空状態における光の伝播法則をより複雑なものに置き換えるかのいずれかを選択せざるを得ないようである。あなたがたが前段落の議論を慎重に追跡した場合、相対性原理を留保すべきだと期待するのはほぼ確実であろう。この原理はいかに自然かつ単純であり、またそのような説得力があるからである。しかし、光の伝播法則を真空状態においてより複雑なものに置き換えることになる。理論物理学の発展は、この道を追うことができないと示している。H.A.ローテンツによる動体に関連する電磁気学的および光学的な現象に関する革新的な研究は、この領域での経験が、相対性原理と光の伝播法則との間に最小限の不一致がないことを結論付ける理論を導くことを見せる。実証データがこの原理に反するものを見つけられなかったにもかかわらず、有名な理論物理学者たちは、相対性原理を放棄すべきであると傾倒した。

この局面においては、相対性理論が舞台に登場した。時間と空間の物理学的概念に対する分析により明らかになったように、実際には相対性原理と光の伝播法則との間に最小限の不一致がないことが明白となった。そして、両方の法則を系統的に堅持することで、論理的な堅固な理論が導き出せることがわかった。この理論は「特別な相対性理論」として呼ばれるもので、後述する拡張された理論と区別するために用いられた。以下のページでは、特別な相対性理論の基本的なアイデアを提示する。


\clearpage
\chapter{第VIII節：On the Idea of Time in Physics}

第八章。

物理学における時間の概念

鉄道敷設線上の二地点$A$と$B$に雷が同時に打ったとする。私はさらに、これらの二つの雷閃は同時に行われたという主張を加える。この声明の意味について尋ねるなら、あなたは「Yes」と明確に答えるだろう。しかし、その意味をより精密に説明するように求められた場合、最初の印象とは異なるほど答えは簡単ではないことが考慮される。

しばらく考えた後、次のような回答が思い浮かぶかもしれない。「この声明の意味は自明であり、追加的な説明は必要ない。もちろん、実際の二つの事件が同時に起こったかどうかを観測によって判断するためには考慮が必要だ」というものである。しかし、私はこの回答に満足することはできない。なぜなら、巧みな考察により雷が$A$と$B$の両地点で常に同時に打つことを発見したと仮定すると、その理論的な結果が実際と一致しているかどうかを検証する必要があるという課題に直面することになるからである。物理学におけるすべての声明において、「同時」の概念が重要な役割を果たす場合も同様の困難さに遭遇する。この概念は物理学者にとって存在しないまで、実際にその条件が満たされているかどうかを確認する可能性があるか否かを決定する方法が必要である。したがって、実験によって$A$と$B$の両地点で打った雷閃が同時に起こったかどうかを判断するための定義が必要となる。この要件が満たされない限り、物理学者（もちろんそれ以外の人も含む）は「同時」という声明に意味を付与できると錯覚することを許すべきではない。（私は読者に、この点で完全に納得するまで進むのを控えていただくことをお願いします。）

あなたは長く考えた後、次のような提案を提出し、同時性を検証する方法を提示しました。鉄道の軌道上を測定して接続線$A_{B}$を測り、距離$A_{B}$の中点Mに観察者を配置します。この観察者は（例：90°の傾斜角度を持つ2つの鏡など）で、同時に場所\\textit{Aと}Bの雷光を視覚的に観察できる装置を提供されます。もし観察者が両方の雷光を同じ時間に感知した場合、それは同時であると判断できます。

私はこの提案が非常に満足していますが、それでも私はその問題が完全に解決されたとは考えられません。なぜなら、次のような反論が必要不可欠だと感じています。「あなたの定義は確かに正しいかもしれませんが、観察者$M$が雷光を感知するための光が$A$→$M$と$B$→$M$の距離を同じ速度で進むことを知っている限りです。しかし、この仮定の検証は時間測定の手段が既に手元にある場合のみ可能となります。したがってここでは論理的な円周運動をしているように思われます」

さらに考えた後、あなたは少し軽蔑的についには目を向け、「しかし私は以前の定義を維持します」と宣言しました。そして「実際には、その定義は光の物理的な性質について全く仮定していないのです。 同時性の定義には一つだけ要求があります：具体的な事例において、その定義が満たされているかどうかの実証的な判断を提供する必要があります。私の定義がこの要求を満たしていることは争いの余地がありません。光が$A$→$M$と$B$→$M$の両方の経路を同じ時間で進むことについては、実際には光の物理的な性質についての仮定や推測ではなく、自由意志に基づいて 同時性の定義を作成するための規定であると理解できます。」

この定義は、二つの事象だけでなく、任意の数の事象にも正確な意味を付与できることが明らかです。また、参照体（ここでは鉄道線路）に対する事象の場所の位置関係に依存することなく使用できます。これにより物理学における「時間」の定義も導き出されます。この目的のために、同一構造の時計を鉄道路線上の点 $A, B$ および $C$ に設置し、その指針の位置が同時に（上記の意味で）同じであるように調整します。このような条件下では、事象が発生した直近（空間的に）にあるこれらの時計のうちの一つの読み取り（指針の位置）を「時間」と定義します。これにより、観察可能な要素に固有の時間値が付与されます。

さらに、異なる場所で発生する三つの事象 $A, B$ および $C$ について、 $A$ が $B$ と同時であり、 $B$ が $C$ と同時に、（上記定義による）同時性を満たす場合、 $A$ と $C$ の同時性も満足されるという基準があります。この仮定は光の伝播法則に関する物理学的仮説であり、真空中的な光速の恒常性を維持するためには必須です。

この規定にはさらに物理的な仮定が含まれており、その妥当性は実験的証拠がない限りほとんど疑われないでしょう。すべてのこれらの時計が同一構造であれば同じ速度で進むという仮定があります。より正確に述べれば、参照体上の異なる場所にある二つの静止配置された時計を調整し、一つの時計の指針の特定位置が（上記の意味で）別の時計の指針の同一位置と同時に存在する場合、同じ「設定」は常に同時に（上記定義による）存在します。


\clearpage
\chapter{第IX節：The Relativity of Simultaneity}

IX.

同時性の相対性

これまでまで、我々の考慮は特定の参照体に引き継がれてきた。私たちはこれを「鉄道仮設」だと呼んだ。「鉄道仮設」という非常に長い列車が図1で示す方向と速度vで移動していると仮定する。この列車に乗っている人々は、列車を剛体の参照体（座標系）として見なし、列車上でのすべての事象について列車を基準とする。

すると、線路上で起こるすべての事象も列車上の特定の点で起こると考えられる。また、同時性の定義は、列車に対して鉄道仮設に対して同様に与えることができる。

しかし、自然な結果として、以下の質問が生じる：

二つの事象（例えば、光閃AとB）が鉄道仮設に対する 同時性 である場合、それらの事象は列車に対する 同時性 でもあるか？

直接証明により、この問いに対する答えは否定的であることを示すことができる。

当たる瞬間$A$と$B$が列車の静止地における同時性であるとは、雷が打った場所$A$と$B$で発せられた光線が列車の長さ$A$→$B$の中点$M$で交わることを意味します。しかし、事件$A$と$B$は列車上の位置$A$と$B$に対応しています。列車上での距離$A$→$B$の中点を$M′$とします。雷鳴が静止地から見えた瞬間、この点$M′$は自然に$M$と一致しますが、図では右方向に速度$v$で移動します。列車上の位置$M′$に座っている観測者がこの速度を持っていなければ、彼は永久に$M$に留まり続け、雷鳴$A$と$B$から発せられた光線は同時に到達します、すなわち、彼がいる場所で交わるでしょう。実際には（鉄道列車の参照体に関して考慮すると）、彼は$B$からの光束よりも早く$A$からの光束に近づいています。したがって、観測者は$B$より$A$の方が早く光束を見ることになります。つまり、列車を参照体とする観測者は雷鳴$B$が雷鳴$A$より早かったと結論付けるでしょう。これにより重要な結果を得ます：

静止地における同時性は列車に対する同時性とは異なり、逆もまた然り（ 同時性の相対性）。各参照体（座標系）にはそれぞれ固有の時間があります。参考体を明示しない限り、「事件の時刻」の声明には意味がありません。

理論的相対性以前に物理学では、時刻の記述は絶対的な意味を持つと無意識的に仮定されていました、すなわち、参照体の運動状態によらず独立していると考えられてきました。しかし、この仮定は自然な同時性の定義と不一致です。この仮定を捨てれば、真空での光の伝播法則と相対性原理（第7節で開発された）との間の矛盾が解消します。

われわれは第六節の考慮によりその対立に至ったが、現在その節の結論は成り立たない。そこで我々は車内を移動する男が、一秒ごとに車両に対して距離 $w$ を通過すると結論付けた。しかし、前節の考慮によれば、特定の出来事に対する車両での時間は、それと同じ出来事が鉄道線路（参考体）からの観測でどのくらいの期間かを考慮する上で等しくないものとされるべきである。したがって、車内を歩行する男が一秒間として鉄道線路から判断すると同じ時間に距離 $w$ を通過しているとは主張できない。

さらに第六節の考慮は、相対性理論導入以前から無意識的に仮定されていたもう一つの前提に基づいている。この前提は、厳密な検討において任意であると見做されるべきであり、相対性理論が導入された後も常に無言で仮定されていた。


\clearpage
\chapter{第X節：On the Relativity of the Conception of Distance}

X.

距離の概念に関する相対性について

列車が運転されている軌道上にある二つの特定の点を考えましょう。これらの点間の距離はどの程度かを調査します。距離を測定するためには、その距離を測る基準となる物体が必要です。最も単純な方法として、列車自体を参照物（座標系）として使用することができます。列車にいる観察者は、 carriageの床に沿って測量棒を一直線に並べて、一つの点からもう一つの点まで移動する回数を数えます。その数が求められる距離となります。

 例えば、最初の車と百番目の車の中央点。

しかし、列車からの距離を軌道上から判断する場合、状況は異なります。このとき、次のような方法を考えることができます。二つの点 $A′$ および $B′$ の間の距離を求めたい場合、これらは特定の時間 t（軌道上から観測）においてそれぞれ軌道上の点 $A$ と $B$ を通過しています。これらの点 $A$ および $B$ を定めるために、第 VIII 節で与えられた時間の定義を使用することができます。その後、列車上の測量棒を軌道上に沿って繰り返し使用して、点 A と B の間の距離を測定します。

最後の測定が最初の測定と同じ結果を提供するとは必ずしも保証されていません。つまり、列車からの位置で測った列車の長さは、列車内で測った長さとは異なる可能性があります。この事実は、第 VI 節で明らかにされた考えに対する二番目の反論を提起します。具体的には、列車上の男性が一定時間（列車上から測定）で距離 $w$ を覆う場合、その距離（軌道上から測った場合）は必ずしも $w$ とは限らないということです。


\clearpage
\chapter{第XI節：The Lorentz Transformation}

最後の三節の結果は、光の伝播法則と相対性原理（第VII節）との一見した不一致が、古典力学から借用した二つの根拠のない仮説に基づいて導き出されたものであることを示している。それらの仮説とは以下の通りである：

(1) 二つの事象の間の時間間隔（時間）は、基準体の運動状態に依存しない。

(2) 剛体上の二点間の空間的間隔（距離）は、基準体の運動状態に依存しない。

もしこれらの仮説を破棄すれば、第VII節におけるジレンマは消失する。なぜなら第VI節で導出された速度の合成則が無効となるからである。真空中の光伝播法則が相対性原理と両立する可能性が開かれ、そこで次のような問いが生じる。これら二つの基本的な経験的結果の間の見かけ上の不一致を取り除くためには、第VI節の考察をどのように修正すべきであろうか？

この問いはより一般的な問いへと導く。第VI節の議論において、我々は列車と線路堤の両方に対する場所と時間を取り扱った。ある事象の鉄道線路堤に対する場所と時間が分かっているとき、その事象の列車に対する場所と時間をいかにして求めるべきか？ 真空中の光伝播法則が相対性原理に矛盾しないような、考え得る解答は存在するであろうか？ 換言すれば、個々の事象の両基準体に対する場所と時間の間に、あらゆる光線が線路堤に対しても列車に対しても同一の伝播速度 $c$ をもつような関係を構想できるであろうか？

この問いは、極めて明確な肯定的解答へと導かれ、ある基準体から別の基準体へと移行する際の事象の時空諸量に関する完全に決定された変換則をもたらす。

これに入る前に、以下の付随的考察を導入しておこう。これまでのところ、我々は数学的に直線の役割を担うべき線路堤に沿って生じる事象のみを考えてきた。第II節で示した方法と同様に、我々はこの基準体が棒の骨組みによって側面および垂直方向に拡張されていると想像することができ、それによりどこで発生する事象であってもこの骨組みを基準として位置づけることができる。同様に、速度 $v$ で走る列車が全空間にわたって拡張されていると想像することもでき、どれほど遠方で生じる事象であってもこの第二の骨組みに対して位置を特定できる。これらの骨組みが実際には固体の不入性ゆえに互いに衝突・干渉し合うという事実は、根本的な誤りを生じさせるものではない。

このような各々の骨組みにおいて、互いに直交する三つの平面を想定し、これらを「座標平面」（「座標系」）と呼ぶことにする。座標系 $K$ は線路堤に対応し、座標系 $K'$ は列車に対応する。どこで発生した事象であっても、座標平面への三本の垂線 $x, y, z$ によって空間的に固定され、時刻値 $t$ によって時間的に固定される。$K'$ に対しても、同一の事象は対応する値 $x', y', z', t'$ によって時空的に固定されるが、これらは当然ながら $x, y, z, t$ とは一致しない。これらの諸量が物理的測定の結果としてどのように見なされるべきかについては、すでに詳細に述べておいた。



\begin{center}
\begin{tikzpicture}[scale=0.8, >=stealth, thick]
  % System K
  \draw[->] (0,0) -- (3,0) node[right] {$x$};
  \draw[->] (0,0) -- (0,3) node[above] {$y$};
  \draw[->] (0,0) -- (-1.5,-1.5) node[below left] {$z$};
  \node at (-0.3,0.3) {$K$};
  % System K'
  \draw[->, shift={(4.5,0)}] (0,0) -- (3,0) node[right] {$x'$};
  \draw[->, shift={(4.5,0)}] (0,0) -- (0,3) node[above] {$y'$};
  \draw[->, shift={(4.5,0)}] (0,0) -- (-1.5,-1.5) node[below left] {$z'$};
  \draw[->, very thick, blue] (5.2, 1.5) -- (6.5, 1.5) node[right] {$v$};
  \node at (4.2,0.3) {$K'$};
\end{tikzpicture}\\[1mm]
{\small 図 2：互いに平行な空間配置を持つ二つの慣性座標系 $K$ と $K'$}
\end{center}



明らかに、我々の問題は次のように正確に定式化できる。基準体 $K$ に対する同一の事象の諸量 $x, y, z, t$ が与えられたとき、基準体 $K'$ に対するその事象の値 $x', y', z', t'$ は何か？ これらの関係は、真空中における光の伝播法則が、$K$ および $K'$ に対する同一の光線（そしてもちろんあらゆる光線）について満たされるように選ばれなければならない。図2に示された座標系の空間における相対的配置に対して、この問題は以下の方程式系によって解決される：



\[ x' = \frac{x - vt}{\sqrt{1 - \frac{v^2}{c^2}}} \]



$y' = y$

$z' = z$



\[ t' = \frac{t - \frac{v}{c^2}x}{\sqrt{1 - \frac{v^2}{c^2}}} \]



この方程式系は「ローレンツ変換」として知られている。

 ローレンツ変換の平易な導出は付録Iに与えられている。

もし光の伝播法則の代わりに、時間と長さの絶対的性質に関する旧来の力学の暗黙の仮定を基礎としていたならば、上記の代わりに次の方程式系が得られていたはずである：

$x' = x - vt$

$y' = y$

$z' = z$

$t' = t$

この方程式系はしばしば「ガリレイ変換」と呼ばれる。ガリレイ変換は、ローレンツ変換において光速度 $c$ を無限大の値に置き換えることによって得られる。

以下の説明の助けを借りれば、ローレンツ変換に従うとき、真空中の光伝播法則が基準体 $K$ に対しても基準体 $K'$ に対してもともに満たされることを容易に理解できる。光信号が正の $x$ 軸に沿って送られ、この光刺激は方程式

$x = ct$

に従って進む。すなわち速度 $c$ で進む。ローレンツ変換の方程式によれば、$x$ と $t$ の間のこの単純な関係は、$x'$ と $t'$ の間の関係を伴う。実際、ローレンツ変換の最初と第四の式において $x$ に値 $ct$ を代入すると、



\[ x' = ct \sqrt{\frac{1 - \frac{v}{c}}{1 + \frac{v}{c}}}, \quad t' = t \sqrt{\frac{1 - \frac{v}{c}}{1 + \frac{v}{c}}} \]



が得られ、これらを互いに割り算することによって、直ちに

$x' = ct'$

という表現が導かれる。系 $K'$ を基準とするとき、光の伝播はこの方程式に従って生じる。したがって、基準体 $K'$ に対する光の伝播速度もまた $c$ に等しいことが分かる。他のいかなる方向に進む光線についても同様の結果が得られる。これは当然のことであり、驚くべきことではない。なぜならローレンツ変換の方程式は、まさにこの観点に適合するように導出されたものだからである。


\clearpage
\chapter{第XII節：The Behaviour of Measuring-Rods and Clocks in Motion}

XII.

測量棒と時計の運動における挙動

K′座標系の$x′$-軸に1メートル棒を配置し、一方の端（始点）が$x′$ = 0の点と一致し、他方の端（棒の終点）が$x′$ = 1の点と一致させる。この棒の長さは$K$座標系に対する相対的な長さは何であるか？これを知るためには、特定の時間$t$における$K$座標系でのこれらの2点の位置を学ぶ必要がある。ローレンツ変換の第1式を使用すると、時間$t = 0$におけるこれらの2点の値は

\[
x = \frac{1 - v^2/c^2}{1 - v^2/c^2} = 1
\]

\[
x' = \frac{-v}{c^2}
\]

であり、これら2点間の距離は

\[
\sqrt{\left(1 - \frac{v^2}{c^2}\right)}
\]

となる。つまり、棒は長さが

\[
L = \frac{1}{\sqrt{1 - v^2/c^2}}
\]

メートルである。すなわち、棒は運動しているときよりも静止しているときに短くなるし、その速度が速いほど短くなる。速度$v$ = $c$の場合には

\[
L = \frac{1}{\sqrt{1 - 1}} = \infty
\]

となり、より大きな速度では平方根が虚数になる。これから、相対性理論において速度$c$は制限速度であり、どんな実体もこの速度を達成したり超えたりすることはできないという結論を得る。

もちろん、速度$c$が制限速度であるという特徴もローレンツ変換の式から明らかに導け、これらの式が$v$ > $c$となる値を選んだ場合、それらは意味をなさないからである。

逆に、K座標系の$x$-軸上に静止している1メートル棒を考えた場合、その長さが$K′$からの観測では

\[
L' = \frac{1}{\sqrt{1 - v^2/c^2}}
\]

となることがわかる。これは相対性原理に基づく考察の基礎となっている。

事前に明らかに、変換方程式から測量棒と時計の物理的な挙動について何か学べるということが予想でき、$z, y, x, t$は測量棒や時計を使用して得られる測定結果であることを理解している。我々がガリレオ変換に基づいて考察を行った場合、その運動によって棒の長さが縮小されることを得ることはなかった。

現在、原点（$x′$ = 0）に永久的に設置された秒単位の時計について考えます。$t′$ = 0 と $t′$ = 1 はこの時計の二つの連続する刻みです。ローレンツ変換の最初と第四の方程式により、これらの二つの刻みに対して以下の式が得られます：

$t = 0$

そして

\[
\frac{1}{\sqrt{1 - \frac{v^2}{c^2}}}
\]


$K$ から見れば、時計は速度 $v$ で動いているため、この参照系からは時計の二つの刻み間が経過する時間は 1 秒ではなく、

\[
\frac{1}{\sqrt{1 - \frac{v^2}{c^2}}}
\]


秒、すなわちやや長い時間が経過します。時計の運動により、その動きは静止状態よりも遅くなります。また、ここでも速度 $c$ は到達不能とする極限速度の役割を果たします。


\clearpage
\chapter{第XIII節：Theorem of the Addition of Velocities. The Experiment of Fizeau}

十三章。

速度の加法定理。フィゼーの実験

実際には、我々は光速と比較して小さな速度で時計や測定棒を動かすことができるだけであるから、前節の結果を直接現実と比較することはほとんど不可能であろう。しかし、他方でこれらの結果は非常に不思議に映るだろうから、理論からもう一つの結論を導き出し、その結論は前節の考察から容易に導くことができ、そしてそれは最上級の実験によって確認されている。

第六章では、古典力学の仮定からも得られる形で、一方の方向での速度の加法定理を得た。この定理は、第十一章のガリレオ変換からも容易に導くことができる。車の中で歩いている男を点として置き換え、座標系 \( K' \) において相対的に動く点について、

\[ x' = wt' \]

という方程式を考える。

ガリレオ変換の最初と第四の方程式を使って \( x' \) と \( t' \) を \( x \) と \( t \) の関数として表すことができ、それにより

\[ x = (v + w)t \]

を得る。この式は点が \( K \) 系（または岸に対する男）に参照して運動する法則を表現している。この速度を符号 \( W \) で表し、第六章と同様にして、

\[ W = v + w \] . . . . . . . (A)

を得る。

しかし、相対性理論に基づいても同じような考慮を行うことができる。方程式

\[ x' = wt' \]

において \( x' \) と \( t' \) を \( x \) と \( t \) の関数として表すために、ローレンツ変換の最初と第四の方程式を使う。そうすると方程式 (A) は次の式に置き換わる。

\[ W = \frac{v + w}{\sqrt{1 - \frac{(v + w)^2}{c^2}}} \] . . . . . . . (B)

ここで \( c \) は光速である。

これは相対性理論に基づく一方の方向における速度加法定理に対応します。現在、これらの2つの定理のうち、どの方が経験とより一致しているかという問題が生じます。この点について、半世紀以上前に優れた物理学者フェジューによって行われた重要な実験により明確になります。その後、いくつかの最上級の実験物理学家中の誰もがその結果を確認しており、疑問は存在しません。

この実験は次の問いに関連しています。静止した液体中では特定の速度 $w$ で光が伝播します。それに対して、管 $T$ の上部に流れる液体の速度が $v$ の場合、光は管の方向でどのくらい速さで伝播するでしょうか？

図3を参照してください。

相対性原理に基づき、我々は当然、液体が他の物体に対する運動状態に関係なく、その伝播速度は常に液体に対して一定の速度 $w$ で光が伝播することを仮定しなければなりません。液体に対する光の速度と管に対する液体の速度は既知であり、求められるのは管に対する光の速度です。

これは第6節の問題と同じものであることが明らかです。管は鉄道の線路や座標系 $K$ の役割を果たし、液体は車両や座標系 $K′$ の役割を果たし、最後に光は車両上を歩く人や現在のセクションにおける動点の役割を果たします。光が管に対する速度 $W$ である場合、これはガリレイ変換またはローレンツ変換に対応する方程式（A）または（B）によって与えられます。実験は、ローレンツ変換から導かれた式（B）を支持し、その一致は非常に正確です。最近で最も優れた測定によりゼーマンが行った測定によれば、流速 $v$ の影響は1\%未満の精度で式（B）によって表されます。

  フェジューは以下を発見しました

\[
w' = \frac{w}{1 + \frac{vw}{c^2}}
\]

ここで、\[
n = \frac{c}{w}
\] は液体の折射率です。一方で、\[
\frac{vw}{c^2}
\] が1に比べて非常に小さいため、式（B）はまず以下のように近似できます。

\[
w' \approx w - v
\]

または同程度の近似では、

\[
w' \approx \frac{wv}{c^2}
\]

これはフィザウの結果と一致する。

しかし、理論相対性説が提唱される以前に、この現象について電磁気動力学的な理論を提示した人物はH.A.ローテンスであることに注意しなければならない。この理論は純粋に電磁気動力学的であり、物質の電磁気構造に関する特定の仮説を使用して得られたものである。しかし、この事実は、相対性原理の理論がこの実験による決定的な検証においてその結論を損なうものではない。マクスウェル-ローテンスの電磁気動力学に基づいていた元の理論に対して、相対性原理は驚くほど単純で一般的な組み合わせと一般化として電磁気動力学の仮説を発展させたものである。これらの仮説は以前は独立して電磁気動力学が構築されたものであった。


\clearpage
\chapter{第XIV節：The Heuristic Value of the Theory of Relativity}

XIV.

相対性理論の推敲的价値

前頁での議論は、以下のようになる。経験から、一方で相対性原理が成り立つことを確信し、他方で真空において光の伝播速度は定数 $c$ とみなす必要があるという結論を得た。これらの二つの仮定を組み合わせることにより、自然現象を構成する事象の座標 $x, y, z$ および時間 $t$ の変換法則を得ることができた。この過程においてはギャリレオ変換ではなく、古典力学とは異なる $ローレンツ変換$ を得た。

光の伝播法則が実際的な知識に基づいて正当化され、その重要な役割を果たしたことは明らかである。しかしローレンツ変換を得た後では、この変換と相対性原理を組み合わせることにより、理論を以下のようにまとめることができる：

自然現象の一般法則は、元の座標系 $K$ の空間時間変数 $x, y, z, t$ ではなく、新たな座標系 $K′$ の空間時間変数 \\textit{x′, y′, z′, t′} を用いたときも、同じ形式の法則に変換されることで構成されるべきである。この関係はローレンツ変換によって与えられる。または簡略化して言えば：自然現象の一般法則はローレンツ変換に関して共変である。

これは相対性理論が自然法則に対して要求する明確な数学的条件であり、この条件を満たすための理論は、自然法則探求において有用な推敲手段となる。もし一般法則がこの条件を満たさないものであれば、少なくとも二つの基本的な仮定の一つは否定されることになる。では、この理論がこれまでに示した一般的な結果を見ていきましょう。


\clearpage
\chapter{第XV節：General Results of the Theory}

これまでの考察から明らかなように、（特殊）相対性理論は電磁気学および光学から生まれ育ったものである。これらの領域において、相対性理論は理論の予測を著しく変えるものではなかったが、理論的構造、すなわち法則の導出を極めて簡潔にし、そして——比較にならないほど重要なこととして——理論の基礎をなす独立な仮説の数を大幅に削減した。特殊相対性理論はマクスウェル＝ローレンツの理論を極めてもっともらしいものとしたため、仮に実験がそれほど決定的に有利な結果を示さなかったとしても、物理学者たちによって広く受け入れられたことであろう。

古典力学は、特殊相対性理論の要請と調和させるために修正を余儀なくされた。しかし、この修正の大部分は、物質の速度 $v$ が光速度に比べて極めて小さくはない急速な運動に対する法則にのみ影響を与える。我々がそのような高速運動を経験するのは電子やイオンの場合に限られ、その他の運動については古典力学の法則からの偏差が微小であるため、実際上表面化することはない。天体の運動については、一般相対性理論を論じるまで考察を差し控えることにする。

相対性理論によれば、質量 $m$ の質点の運動エネルギーは、もはや周知の表現



\[ \frac{1}{2} m v^2 \]



によっては与えられず、次の表現によって与えられる：



\[ \frac{m c^2}{\sqrt{1 - \frac{v^2}{c^2}}} \]



この表現は、速度 $v$ が光速度 $c$ に近づくにつれて無限大へと発散する。したがって、加速を生み出すためにいかに大きなエネルギーが投入されたとしても、物体の速度は常に $c$ 未満にとどまらなければならない。もし運動エネルギーの表現を級数の形で展開すれば、次の式が得られる：



\[ m c^2 + \frac{1}{2} m v^2 + \frac{3}{8} m \frac{v^4}{c^2} + \dots \]





\[ \frac{v}{c} \]



が $1$ に比べて小さいとき、これらの項のうち第三の項は、古典力学において唯一考慮されていた第二の項に比べて常に小さい。第一の項 $mc^2$ は速度を含んでおらず、質点のエネルギーがいかに速度に依存するかという問題のみを扱う場合には考慮する必要がない。その本質的な意義については後ほど述べる。

特殊相対性理論が導いた一般的性質をもつ最も重要な結果は、質量の概念に関するものである。相対性理論が登場する以前、物理学は極めて重要な二つの保存則、すなわちエネルギー保存の法則と質量保存の法則を認めていた。これら二つの基本法則は互いに完全に独立しているように見えた。しかし相対性理論によって、これらは一つの法則へと統合されたのである。我々は今、この統合がいかにしてもたらされたか、そしてそれにどのような意味が付与されるべきかを簡潔に考察しよう。

相対性原理は、エネルギー保存の法則が基準座標系 $K$ に対して成り立つのみならず、$K$ に対して等速直線並進運動を行っているあらゆる座標系 $K'$（要するにあらゆる「ガリレイ」座標系）に対しても成り立つことを要求する。古典力学とは対照的に、このようなある系から別の系への移行においては、ローレンツ変換が決定的要因となる。

これらの前提とマクスウェル電磁気学の基本方程式を組み合わせた比較的単純な考察により、我々は次の結論を導き出すことができる。速度 $v$ で運動している物体が、その過程で速度の変化を被ることなく、放射の形で量 $E_0$ のエネルギーを吸収するとき、その結果として物体のエネルギーは次の量だけ増加する：



\[ \frac{E_0}{\sqrt{1 - \frac{v^2}{c^2}}} \]



 $E_0$ は物体とともに運動する座標系から判定された吸収エネルギーである。

物体の運動エネルギーについて上で与えられた表現を考慮すると、物体の全エネルギーは次のように求まる：



\[ \frac{(m + \frac{E_0}{c^2}) c^2}{\sqrt{1 - \frac{v^2}{c^2}}} \]



したがって物体は、速度 $v$ で運動する質量



\[ m + \frac{E_0}{c^2} \]



の物体と全く同一のエネルギーをもつことになる。それゆえ我々は次のように言うことができる。もし物体が量 $E_0$ のエネルギーを吸収するならば、その慣性質量は次の量だけ増加する：



\[ \frac{E_0}{c^2} \]



物体の慣性質量は定数ではなく、物体のエネルギーの変化に応じて変化する。物体の系の慣性質量は、その系のエネルギーの尺度とさえ見なすことができる。系の質量保存の法則はエネルギー保存の法則と同一のものとなり、その系がエネルギーを吸収も放出もしない場合にのみ成立する。エネルギーの表現を



\[ (m + \frac{E}{c^2}) c^2 \]



の形で書き表すと、これまで我々の注目を惹いてきた項 $mc^2$ は、物体がエネルギー $E_0$ を吸収する前に既に保持していたエネルギーにほかならないことが分かる。

 物体とともに運動する座標系から判定されたもの。

この関係を実験と直接比較することは、現時点（1920年；[注]参照）では不可能である。なぜなら我々が系に与え得るエネルギー変化 $E_0$ は、系の慣性質量の変化として知覚できるほど十分には大きくないからである。



\[ \frac{E_0}{c^2} \]



は、エネルギー変化の前に存在していた質量 $m$ に比べてあまりにも小さすぎる。古典力学が質量保存の法則を独立した妥当性をもつ法則として首尾よく確立できたのは、まさにこの事情によるものである。

[注] 方程式 $E = mc^2$ は、その後今日に至るまで幾度となく徹底的に実証されている。

最後に、根本的な性質をもつ注意を一つ付け加えておきたい。遠隔電磁作用に関するファラデーとマクスウェルの解釈の成功により、物理学者たちは、ニュートンの万有引力の法則のような（媒介物を伴わない）瞬間的な遠隔作用というものは存在しないという確信を抱くに至った。

相対性理論によれば、瞬間的な遠隔作用や無限大の伝播速度をもつ遠隔作用の代わりに、常に光速度による遠隔作用が現れる。このことは、速度 $c$ が本理論において根本的な役割を果たしているという事実と密接に関連している。第II部において、我々はこの結果が一般相対性理論においてどのように修正されるかを見ることになる。


\clearpage
\chapter{第XVI節：Experience and the Special Theory of Relativity}

XVI.

経験と特殊相対性理論

特殊相対性理論がどの程度経験によって支持されているかは、フIZEAUの基本実験に関連して既に述べた理由により、容易には答えられません。特殊相対性理論はマクスウェル-ロレンツ電磁気現象論から結晶化したものです。したがって、電磁気現象論を支持するすべての経験的事実も相対性理論を支持します。特に重要であると考えられるのは、相対性理論が我々に固定星からの光が到達する効果を予測することです。これらの結果は非常に単純な方法で得られ、地球が太陽を中心に回転する際に固定星の観測位置が年間で移動する（ aberration ）という現象や、固定星に対する地球の相対的な運動の径向成分が我々に到達する光の色に及ぼす影響です。後者の効果は、固定星からの光を伝える同一のスペクトル線が地上光源によって生成される場合と比較してわずかにスペクトル位置がずれるという現象を示します（ドッペル原理）。マクスウェル-ロレンツ理論を支持する実験的根拠は、同時に相対性理論を支持するものであり、これらはここで述べる余地はありません。実際には、これらの根拠は理論的な可能性を制限し、他の理論よりも経験によって検証されるときに対抗できる唯一の理論としてマクスウェルとロレンツの理論が存在します。

しかし、これまで得られた二つの実験的事実のクラスは、マクスウェル-ロレンツ理論において補助仮定を導入せざるを得ないものであり、その自体では—すなわち相対性理論を使わない限り—余計なものとして見做されるものです。

知られているように、ケッカーライと放射性物質から放出されるβ線は、非常に軽量で高速な負電荷を持つ粒子（電子）由来である。これらの粒子の運動法則を電磁場の影響下での偏折により精密に研究することができる。

これらの電子についての理論的な扱いでは、自己電動論自体がその性質を説明することができないという困難に直面している。電荷符号が同じである電荷は互いに斥力を持つため、電子由来の負電荷粒子は互いの斥力の影響によって散乱されてしまう。そのためには、これまで私たちには不明な種類の別の力が作用していると考えられる。もし現在、電子の構成する電荷量間の相対的な距離が電子の運動中に変化しない（古典力学における剛性接続）と仮定すると、その運動法則は経験とは一致しないものとなる。純粋に形式的な観点から導かれたH.A.ロレンツは最初、運動により電子の形が運動方向に対して収縮するという仮説を導入した。この収縮長さは以下の式で比例している。

\[
L = L_0 \sqrt{1 - \frac{v^2}{c^2}}
\]

この仮説は電動力学的事実に基づいていないが、近年精密に確認されている特定の運動法則を提供する。

一般相対性理論により、電子の電荷量は重力による相互作用によって保持される可能性がある。

相対性理論は、電子の構造や行動に関する特別な仮説を必要とせずに同じ運動法則を得ることができる。第13節でフィザウの実験に関連して結論付けたように、相対性理論は液体の物理的な性質に依存することなくその結果を予測することができる。

この論文で言及した第二の事実は、地球が空間において運動しているかどうかの運動を地上の実験で感知できるかどうかに関する質問に関連しています。第五節で既に述べたように、このような試みはすべて否定的な結果を導き出しました。相対性理論が提案される以前には、この否定的な結果を受け入れるのは困難でした。その理由は後に説明します。時間と空間に対する遺伝的偏見は、ガリレオ変換が異なる参照系に切り替える際に極めて重要であるという疑念を生じさせませんでした。マクスウェル・ロレンツ方程式が参照体 $K$ において成り立つ場合、参照体 $K′$ が $K$ に対して一様に運動していると仮定すると、ガリレオ変換の関係が座標系 $K$ と $K′$ の間に存在すると仮定したとき、マクスウェル・ロレンツ方程式は $K′$ において成り立たないことがわかります。これにより、すべてのガリレオ座標系の中で、特定の運動状態に対応する一つ（$K$）が物理的に一意であることが明らかになります。この結果は、$K$ を仮想的な空間エーテルに対して静止していると解釈して物理的に捉えました。一方で、$K$ に対する相対運動のすべての座標系 $K′$ はエーテルに対して動いているとみなされました。$K′$ のエーテルに対する運動（「エーテル流速」）に、$K′$ において想定されるより複雑な法則が付与されました。厳密には、地球表面でのエーテル流速の存在も仮定すべきであり、長らく物理学者たちは地球表面でのエーテル流速を検出する試みに力を注いできました。

これらの試みのうち最も有名なものの一つで、ミチエルソンは決定的な方法を考案しました。二つの鏡を剛体上に配置し、反射面が互いに向かい合っていると想像してください。光線が一方の鏡から他方の鏡へ往復するには、全体がエーテルに対して静止している場合、一定の時間 $T$ が必要です。しかし、計算によれば、体と鏡がエーテルに対して相対的に動いている場合、この過程のために必要な時間 $T′$ はやや異なります。

さらに重要な点：計算によって示されるように、特定のエーテルに対する速度 $v$ に対して、鏡の平面に垂直に動いているときと平行に動いているときに得られる時間 $T′$ は異なる値になります。これらの二つの時間の差は非常に小さいものの、ミチエルソンとモーリーは、この差が明確に検出できる実験を行いました。しかし、その実験結果は否定的であり、物理学者にとって困惑するものでした。

ロレンツとフィッツジェラルドは、体のエーテルに対する運動によって体が動く方向に収縮すると仮定して理論から難題を解消しました。この収縮量は、先述した時間差を補償するために適切な大きさです。第12節との比較により、相対性理論の立場からはこの解決策が正しいことが示されます。しかし、相対性理論に基づく解釈は圧倒的に満足感があります。この理論によれば、特別に選ばれた（唯一の）座標系を導入するためにエーテルの概念を必要とするものではなく、したがってエーテルの流れやその証明実験はありません。ここでは、運動体の収縮は基本原理から直接導き出され、特別な仮定を導入することなく得られます。そしてこの収縮の主な要因は、自己に意味を与えることができない運動自体ではなく、特定のケースに対する参照系を選んだときのその運動に関するものとなります。したがって、地球とともに動く座標系ではミチエルソンとモーリーの鏡システムは短くなりませんが、太陽に対して静止している座標系では鏡システムは短くなります。


\clearpage
\chapter{第XVII節：Minkowski's Four-Dimensional Space}

非数学者は、「四次元」という言葉を聞くと、まるで心霊術やオカルティズムの思想によって呼び起こされるような感情、すなわち神秘的な戦慄に囚われる。しかし、我々が生きているこの世界が四次元の時空連続体であるという言明ほど、平凡で当たり前の主張は存在しない。

空間は三次元の連続体である。このことが意味するのは、（静止した）点の位置を三つの数値（座標）$x, y, z$ によって記述できるということであり、かつこの点の近傍には無数の点が存在して、それらの位置は、最初の点の座標 $x, y, z$ に我々が望むだけ近い座標 $x_1, y_1, z_1$ によって記述できるということである。この後者の性質ゆえに我々はこれを「連続体」と呼び、三つの座標が存在するという事実ゆえにこれを「三次元」と呼ぶのである。

同様に、ミンコフスキーによって端的に「世界」と呼ばれた物理現象の世界も、時空的な意味において当然ながら四次元である。なぜならそれは個々の事象から構成されており、それらの各事象は四つの数値、すなわち三つの空間座標 $x, y, z$ と、時間座標である時刻値 $t$ によって記述されるからである。「世界」はこの意味においてやはり連続体である。なぜなら、あらゆる事象に対して、最初に考慮された事象 $x, y, z, t$ の座標から無限小の差しかない座標 $x_1, y_1, z_1, t_1$ をもつ「近傍の」事象が（現実に生じるか、あるいは少なくとも考え得るものとして）我々が選ぶだけ無数に存在するからである。我々が世界をこのような意味での四次元連続体として見なすことに慣れていなかったのは、相対性理論が登場する以前の物理学においては、時間が空間座標と比較して異質な、より独立した役割を果たしていたという事実に起因する。まさにこの理由から、我々は時間を独立した連続体として扱う習慣をつけてきたのである。実際、古典力学によれば時間は絶対的であり、すなわち座標系の位置や運動状態から独立している。このことはガリレイ変換の最後の式（$t' = t$）に明確に表れている。

相対性理論において「世界」を四次元的に考察する方法が自然であるのは、この理論によれば時間がその独立性を奪われるからである。このことはローレンツ変換の第四式によって示されている：



\[ \Delta t' = \frac{- \frac{v}{c^2} \Delta x}{\sqrt{1 - \frac{v^2}{c^2}}} \]



さらに、この方程式によれば、$K$ を基準とした同一の二つの事象の時間差 $\Delta t$ がゼロである場合でも、$K'$ に関するそれらの事象の時間差 $\Delta t'$ は一般にはゼロにならない。$K$ に関する二つの事象の純粋な「空間的距離」が、$K'$ における同一の事象の「時間的間隔」を生み出すのである。

しかしながら、相対性理論の形式的発展にとって極めて重要であったミンコフスキーの発見は、この点にのみあるのではない。むしろそれは、相対性理論の四次元時空連続体が、その最も本質的な形式的特性において、ユークリッド幾何学の三次元空間連続体と極めて顕著な親密な関係を示すという事実を彼が認識した点にある。

この関係を十全に際立たせるためには、我々は通常の時間座標 $t$ を、それに比例する虚数量



\[ x_4 = i c t \quad (i = \sqrt{-1}) \]



によって置き換えなければならない。この条件下において、（特殊）相対性理論の要請を満たす自然法則は数学的形態をとり、その中で時間座標は三つの空間座標と全く同じ役割を果たすことになる。形式的に、これら四つの座標はユークリッド幾何学における三つの空間座標と完全に一致する。我々の知見に対するこの純粋に形式的な補足の結果として、理論が並大抵でない明晰さを獲得したことは、数学の専門家でない者にとっても明白であろう。

 付録IIにおけるやや詳しい議論を参照されたい。

これらの不十分な論述は、ミンコフスキーによってもたらされた重要な着想について、読者に漠然とした観念を与えるに過ぎない。この着想がなければ、本書の以降の頁で展開される基本的思想をもつ一般相対性理論は、おそらく産着を脱ぎ捨てる（揺籃期を脱する）ことはできなかったであろう。ミンコフスキーの著作は数学に不慣れな者にとっては近づきがたいものに違いないが、特殊あるいは一般相対性理論の基本的思想を理解するためにその著作を極めて厳密に把握することは必須ではないため、ここではこの話題を措き、第II部の終わり近くで再び立ち返ることにする。

第II部：一般相対性理論


\clearpage
\part{一般相対性理論（The General Theory of Relativity）}



\clearpage
\chapter{第XVIII節：Special and General Principle of Relativity}

十八章。

特別かつ一般の相対性原理

我々がこれまで考慮した全ての基本的な原則は、特別な相対性原理であり、すなわち、すべての均質運動に対する物理的相対性の原理であった。一度 lagi その意味を慎重に分析してみよう。

常に明確だったのは、我々が持つ概念から見れば、どの動きも相対的な動きとしてのみ考慮されるべきであるということだ。しばしば用いた例え話に戻って、列車と鉄道路線沿いの比較を考えると、ここでの動きは次の二つの形で表現できる：

(a) 列車は鉄道路線沿いに対して動いている,

(b) 鉄道路線沿いは列車に対して動いている。

(a) では鉄道路線沿いが、(b) では列車が参照体として用いられる。運動の検出や記述においては、原則としてどの参照体を用いても問題ない。既に述べた通り、これは自明なことだが、「相対性原理」というより広範な声明と混同してはならない。我々が調査の基礎としたものだ。

我々が用いた原理は、我々が観察から得た一般法則を記述する際も、列車や鉄道路線沿いどちらを選んでも等しくよいというだけでなく（これは自明である）、さらに主張していることは次の通り：自然現象の一般的な法則を経験から導き出す際に、

(a) 鉄道路線沿いを参照体として,

(b) 列車を参照体として,

それぞれ記述した場合、これらの一般法則（例えば力学の法則や真空における光の伝播法則）はどちらの場合も同じ形で表されるべきである。これはまた次のように表現できる：自然現象の物理的な記述において、参照体 $K, K′$ のいずれも特別に選ばれているわけではない（「特徴的に区別されている」という意味）。最初の声明とは異なり、この後者の声明は必ずしも先ずから成り立つものではなく、それは「運動」や「参照体」という概念には含まれていないし、それらから導けないものであり、その正しさや誤りは経験によって決めるべきである。

しかし、これまでにおいても、自然法則の公式化に関連してすべての参照系 $K$ の等価性を保証するわけにはいきませんでした。私たちの道筋は次の通りでした。まず第一に、運動状態がガリレオの法則を満たすような参照系 $K$ が存在することを仮定しました：自由に動く粒子は十分な他の粒子から遠ざかっている限り、一直線上を等速に移動します。参照系 $K$(ガリレオ参照系)に関して自然法則はできるだけ単純でなければなりません。しかし、さらに $K$ に対して \\textit{均匀直線運動と回転なしの運動状態} を持つすべての参照系 $K′$ も優先的に扱われるべきであり、これらの参照系は自然法則を公式化する上で $K$ と完全に等価であるとみなされるべきです。これら全ての参照系はガリレオ参照系として扱われます。相対性原理の有効性はこれらの参照系に対してのみ仮定され、他の参照系(\\textit{例えば} 違った種類の運動を持つもの)に対しては仮定されませんでした。この意味で私たちは特別な相対性原理、または特殊相対性理論と呼びます。

これとは異なり、「一般相対性原理」とは以下の主張を指します：すべての参照系 $K, K′$, などは自然現象(\\textit{つまり} 自然法則)の説明において等価であり、どのような運動状態であってもそのように考えられます。しかし、さらに進む前に、これは後の段階で明らかになる理由により、より抽象的な形に置き換えるべきです。

自特殊相対性原理の導入が正当化された以来、一般化を追求する全ての知性は、一般相対性原理への一歩を踏み出す誘惑を感じるだろう。しかし、単純で表面的にも頼りになる考慮によれば、少なくとも現時点ではこのような試みに成功する希望は薄いように思われる。想像してみよう。我々が古い友人である列車の車内に移動したと仮定し、その車両が一定速度で進んでいるとする。一定速度で進んでいる間は、車内の乗客はその動きを感知しないため、彼は列車が静止していると考え、一方で軌道が動いていると解釈することを選択できる。さらに、特殊相対性原理によれば、この解釈は物理的な観点からも正当化される。

しかし、現在の動きが非一様なものに変更されると、例えばブレーキを強力に使用した場合、車内の乗客はそれに応じて前方へ強い衝撃を感じる。この減速運動は列車内の人々に対する物体の力学的挙動を通じて現れる。その力学的挙動は以前考慮されたケースとは異なるため、非一様に動いている列車に対して静止または一定速度で動いている列車における同じ力学法則が適用されることはないように思われる。少なくとも、ガリレオの法則は非一様に動いている列車に対して適用されないことが明らかである。

この理由から、我々は現在この結論を認めざるを得ず、一般相対性原理と反対に非一様運動には一種の絶対的な物理的現実性があると考える。しかし、その後ではこの結論が維持できないことをすぐに見ることができる。


\clearpage
\chapter{第XIX節：The Gravitational Field}

十九章.

重力場

「石を拾って放すと、なぜそれが地面に落ちるのだろうか？」この質問に対する一般的な答えは：「それは地球によって引き寄せられるからだ」というものである。現代物理学では、次のような理由により、その回答を異なる形で述べることが好まれる。

電磁気現象のより精密な研究により、遠距離での作用を無介質空間を通じて直接的に行うと考えることは不可能だと認識されるようになった。例えば、磁石が鉄片を引き付ける場合、磁石が空洞空間を通じて鉄片に直接的に作用するという解釈は不十分であり、フェイデナーやそのような考え方に従って、磁石は常に周囲の空間で物理的な何かを作り出すと考えられる。それが「磁場」と呼ばれるものである。この磁場が鉄片に作用し、鉄片は磁石に向かって動くよう努力する。

ここでは、この付随的な概念について議論することはしない。その正当性については、確かに多少恣意的である。しかし、この概念を用いることで、電磁気現象を理論的により適切に表現できることが示され、特に電磁波の伝播においてそれが特に重要であることが確認されている。重力作用も同様に扱われる。

地球が石に対して行う作用は間接的である。地球はその周囲で重力場を作り出し、この重力場が石に作用し、落下運動を生じさせる。経験から知るところでは、物体に対する作用の強度は地球からの距離によって一定の法則に従って減少する。私たちの見解からは、これは空間における重力場の性質を制御する法則が完全に定められたものでなければならないことを意味し、それが正確に重力作用の距離による減少を表現できるからである。それは次のようなものである：物体（例えば地球）はその直近の周囲で直接的に場を作り出す；その物体から離れた点での場の強度と方向は、空間における重力場の性質が制御される法則によって決定される。

電磁場とは異なり、重力場には最も特徴的な性質が見られ、これこそ後の議論において重要な意味を持つ。単に重力場の影響を受けている物体は加速を受けるが、その加速は物体の材質や物理的状態によらない。例えば、鉛と木の片方が静止状態から或いは同じ初速度で真空状態における重力場内では同じように落下する。この法則は非常に正確に成り立つものであり、以下の考慮に基づいて異なる形で表現できる。

ニュートンの運動方程式によれば、

（力） = （慣性質量） × （加速度）

である。ここで「慣性質量」は加速体の特性的定数である。もし重力が加速の原因であれば、

（力） = （重力質量） × （重力場の強度）

となる。ここでも「重力質量」は物体に対する特性的定数である。これらの2つの関係から、

\[
\frac{\text{重力質量}}{\text{慣性質量}} = \frac{\text{重力場の強度}}{\text{加速度}}
\]

が導かれる。

経験によれば、加速は物体の種類や状態に依存せず、一定の重力場において常に同一であるとすれば、重力質量に対する慣性質量の比もまた全ての物体に対して同じである。適切な単位を選択することでこの比を1とすることができる。したがって我々には以下の法則がある：

（物体の）重力質量 = （物体の）慣性質量

これは重要な法則であり、これまで機械学において記録されてきたが、その解釈はなかった。適切な解釈を得るためには以下を認識することが必要である：物体の同一特性は状況によって「慣性」または「重力」（「 heaviness」と訳す）として現れる。

次節ではこの事実がどの程度成り立つか、そしてその問題が相対性原理の一般的な公理とどのように関連しているかを示す。


\clearpage
\chapter{第XX節：The Equality of Inertial and Gravitational Mass as an Argument for the General Postulate of Relativity}

我々は、恒星やその他の目立った質量から遠く離れた、広大な真空の宇宙空間を想像してみよう。そこでは、ガリレイの基本法則が要求する条件が近似的に満たされている。このとき、この空間（世界）に対して、静止している点は静止し続け、運動している点は永久に等速直線運動を続けるようなガリレイ基準体を選ぶことが可能である。

この基準体として、観測装置を備えた観測者が中に入った、ひとつの部屋のような広々とした「箱」を想像してみよう。この観測者にとっては、当然ながら重力は存在しない。彼は紐で自らの体を床に固定しておかなければならず、さもなければ床をわずかに蹴っただけでも、部屋の天井に向かってゆっくりと浮き上がってしまうであろう。

いま、箱の蓋（天井）の中央の外側にロープを取り付けたフックを固定し、ある「存在」（どのような存在であるかは我々にとって問うところではない）が、一定の力でこのロープを引っ張り始めたとしよう。すると箱は観測者とともに、「上方」に向かって等加速度運動を始める。時が経つにつれて、彼らの速度は途方もない値に達するであろう——もちろん、ロープで引っ張られていない別の基準体からこのすべてを観察しているとすれば、の話である。

しかし、箱の中の人間はこの過程をどのように受け止めるであろうか？ 箱の加速度は、箱の床からの抗力（反作用）によって彼に伝達される。したがって、床の上に大の字に倒れたくないならば、彼は自らの両足によってこの圧力を支えなければならない。このとき彼は、我々の地球上にある部屋の中に立つ人間と全く同じようにして、箱の中に立っているのである。もし彼が手の中に持っていた物体を放したとすれば、箱の加速度はその物体にはもはや伝達されないため、その物体は加速された相対運動によって箱の床へと近づいていく。観測者はさらに、実験にどのような種類の物体を用いようとも、物体が箱の床へ向かって落下する加速度は常に同一の大きさをもつことを確信するであろう。

前節で議論された重力場に関する知識に依拠して、箱の中の人間は、自分と箱が時間的に一定な重力場の中に置かれているという結論に達する。もちろん彼は、なぜ箱がこの重力場の中で落下しないのか、一瞬不思議に思うであろう。しかし、まさにそのとき、彼は箱の蓋の中央にあるフックとそれに結ばれたロープを発見し、その結果、箱はこの重力場の中に静止して吊り下げられているのだという結論に至るのである。

我々はこの男を愚かだと笑い、その結論が誤っていると言うべきであろうか？ 思考の一貫性を保ちたいと願うならば、決してそうすべきではないと私は考える。むしろ、彼の状況把握の方法は、理性にも、既知の力学法則にも全く反していないことを認めねばならない。箱が最初に考察した「ガリレイ空間」に対して加速されているとしても、我々は依然として箱を静止していると見なすことができるのである。こうして我々は、相対性原理を互いに加速し合う基準体にまで拡張するための確固たる根拠を得たのであり、その結果として、一般化された相対性の要請に対する強力な論拠を獲得したのである。

この解釈の可能性が、すべての物体に同一の加速度を与えるという重力場の根本的特性に、あるいは同じことであるが、「慣性質量と重力質量の等価性の法則」に基づいている点に、我々は細心の注意を払わなければならない。もしこの自然法則が存在しなかったならば、加速される箱の中の人間は、周囲の物体の挙動を重力場の仮定に基づいて解釈することはできず、経験を論拠として自らの基準体が「静止している」と想定することは正当化されなかったであろう。

いま、箱の中の人間が蓋の内側にロープを固定し、ロープの自由端に物体を結びつけたと仮定しよう。その結果としてロープは引き伸ばされ、「鉛直」下方に垂れ下がる。もしロープに生じる張力の原因について尋ねられたなら、箱の中の人間は次のように答えるであろう。「吊るされた物体は重力場の中で下方への力を受けており、これがロープの張力によって打ち消されているのだ。ロープの張力の大きさを決定しているのは、吊るされた物体の『重力質量』である」と。

他方、空間に自由に浮かんでいる観測者は、この事態を次のように解釈するであろう。「ロープは必然的に箱の加速度運動を共に分担しなければならず、その運動を結びつけられた物体へと伝達しているのだ。ロープの張力は、物体を加速させるのにちょうど十分な大きさである。ロープの張力の大きさを決定しているのは、物体の『慣性質量』である」と。

この例に導かれて、我々は相対性原理の拡張が、慣性質量と重力質量の等価性の法則の『必然性』を含意していることを知る。こうして我々はこの法則の物理学的解釈を得たのである。

加速される箱の考察から明らかなように、一般相対性理論は重力の法則に関して重要な結果をもたらすに違いない。実際、相対性の一般的着想を系統的に追究することによって、重力場が満たすべき法則が与えられたのである。しかし、これ以上進む前に、私はこれらの考察から生じかねないひとつの誤解に対して読者に警告しておかねばならない。最初に選ばれた座標系に対しては重力場が存在しなかったという事実にもかかわらず、箱の中の人間にとっては重力場が存在している。そこで我々は、重力場の存在とは常に単なる『見かけ上』のものに過ぎないのではないかと軽率に考えてしまうかもしれない。あるいは、どのような種類の重力場が存在していようとも、それに対して重力場が『存在しない』ような別の基準体を常に選ぶことができると考えるかもしれない。しかし、これは決してすべての重力場について真なのではなく、極めて特別な形式をもつ重力場についてのみ成り立つのである。例えば、地球の（全体の）重力場がそこから見れば完全に消失してしまうような基準体を選ぶことは不可能である。

いまや我々は、第XVIII節の最後において一般相対性原理に対して提起したあの議論が、なぜ説得力を持たないのかを理解することができる。列車の中の観測者がブレーキをかけた結果として前方への衝撃（ガタつき）を経験し、これによって列車の運動の非一様性（減速）を認識することは確かに事実である。しかし彼は、この衝撃を列車の『実在の』加速度（減速）に帰着させることを何人からも強制されない。彼は自らの経験を次のように解釈することもできるのである。「私の基準体（列車）は恒久的に静止し続けている。しかしそれに対して（ブレーキがかけられている間）、前方に向かい、時間とともに変化する重力場が存在しているのだ。この重力場の影響のもとで、線路堤は地球とともに非一様に運動し、後方への本来の速度が連続的に減殺されていくのである」と。


\clearpage
\chapter{第XXI節：In What Respects are the Foundations of Classical Mechanics and of the Special Theory of Relativity Unsatisfactory?}

XXI.

古典力学と特別相対性理論の基礎が不十分である点は何か？

私たちはすでに何度も述べたように、古典力学は以下の法則から出発します：十分に他の物質粒子から遠ざかっている物質粒子は、一直線を継続して均一な運動を続けるか、または静止状態を維持します。また、この基本的な法則が特定の独特な運動状態を持つ参照系 $K$ においてのみ有効であることを強調してきました。これらの参照系 $K$ は互いに等速移動しています。他の参照系 $K$ に対しては、この法則は有効ではありません。古典力学と特別相対性理論では、認識された「自然の法則」が成り立つ参照系 $K$ と、それらの法則が成り立たない参照系 $K$ を区別します。

しかし、論理的な思考を持つ者はこの状況に満足することはできません。彼は問います：「なぜ特定の参照系（またはその運動状態）が他の参照系（またはその運動状態）よりも優遇されるのか？この優遇の理由は何なのか？」この質問の意味を明確にするために、比較を使って説明します。

私はガスレンジ前に立っています。レンジ上には、互いに非常に似た2つの鍋が並んでおり、どちらも半分ほど水で満たされています。一方の鍋から蒸気が継続的に放出されていることに気づきますが、もう片方からは何もありません。これは驚くべきことですが、私はガスレンジや鍋を見たことがないとしても驚きを感じます。しかし、最初の鍋下に青白い光る何かを発見し、二つ目には何も見られない場合、私は驚きは和らぎますが、初めてガス炎を見たことがないとしても驚きは消失しません。なぜなら、青白いものは蒸気の放出を引き起こす原因であると説明できるからです。あるいは、少なくともその可能性があります。しかし、どちらの鍋にも青白いものが見られず、一方だけが継続的に蒸気を放出し続けている場合、私は驚きと不満を感じ続けるまで、二つの鍋の異なる行動に atribuibleな状況を見つけるまでです。

古典力学（または特別相対性理論において）では、異なる参照系$K$と$K′$に対する物体の異なる挙動を説明する実在の何かを探しても無駄であると私は虚しい思いをする。ニュートンはこの反論を見抜き、それを否認しようと試みたが、成功しなかった。一方でエーリヒ・マッハは最も明白にこの反論を指摘し、そのためメカニクスは新しい基礎に置かれるべきだと主張した。その反論から解放される唯一の方法は、一般相対性原理に適合する物理学を通じて得られるものである。このような理論の方程式はどの参照系においても成り立つため、どのような運動状態でも有効である。


\clearpage
\chapter{第XXII節：A Few Inferences from the General Principle of Relativity}

XXII.

一般相対性原理からいくつかの推論

第XX節の考慮により、一般相対性原理は我々に、重力場の性質を純粋理論的手法で導き出す立場を提供します。例えば、私たちは任意の自然過程が、ガリレオ領域においてガリレオ座標系$K$に対する方法で進行する空間時間「経路」を知っていると仮定しましょう。純粋な理論的操作（すなわち単なる計算）により、それによっては、加速された座標系$K′$から見たこの既知の自然過程がどのように見えるかを見つけることができます。しかし、新しい座標系$K′$に対する重力場が存在する場合、私たちの考慮もまた、重力場が研究した過程にどのような影響を与えるかを教えてくれます。

例えば、ガリレオが定めた法則によれば、$K$に対する直線的且つ均一な運動をしている物体は、加速された座標系$K′$(船体)に対して、加速度付で一般には曲線的な運動を実行していることが分かります。この加速度や曲率は、$K$に対する重力場の影響に対応しています。既知の事実として、重力場は物体の運動にこのような影響を与えることを理解しており、したがって私たちの考慮は特に新しいものではありません。

しかし、光線について類似の考慮を進めるときには、新たな基本的な結果を得ることができます。ガリレオ座標系$K$に対する場合、そのような光線は速度$c$で直線的に伝搬します。加速された船体（座標系$K′$）に対する視点からは、同じ光線の経路が一直線ではなくなります。これから、一般に重力場では光線が曲線的に伝播することを結論付けることができます。

この結果は二つの面で非常に重要です。

まず第一に、これは実際の現象と比較できるものである。詳細な検討によれば、一般相対性理論によって求められる光線の曲率は、実際的な我々の手元にある重力場では極めて小さいが、太陽に近接して光線が通過する場合のその曲率の推定値は1.7秒の弧である。これは次のように現れるべきである。地球から見れば、特定の固定星は太陽の近くに見え、日食の際にこれらの星を観測できる。このような時、これらの星は上記の量により太陽から外側にずれ込むべきである。そのずれ込みが天体の位置と比較して現れるべきである。この推論の正しさや誤りを検証する問題は最も重要な問題であり、天文学者による早期解決が期待される。

 これは、ロイヤル・アストロノミカル協会共催委員会によって装備された2つの探査隊の星写真により、1919年5月29日の日食時に理論によって要求される光線の偏角が初めて確認された。 （参照：付録III）

次に、我々の結果は、一般相対性理論によれば、特殊相対性理論の二つの基本的な仮定の一つである真空での光速の恒常性の法則（その他の一つは空間の局所性）が無制限の有効性を主張することはできないことを示している。光線の曲率が発生するには、光の伝播速度が位置によって変化することが必要である。これにより、特殊相対性理論とそれに伴う全体的な相対性理論が打ち倒されるべきだと考えるかもしれないが、実際はそうではない。我々は結論を出すことができる。特殊相対性理論は無制限の有効性を主張することはできない。その結果は、重力場の影響を考慮に入れることが可能である限り（例えば光の場合）にのみ成り立つ。

相対性理論の反対派が特別相対性理論を一般相対性理論によって覆われたと主張することがしばしばあるので、適切な比較によって事実を明確にするのは適切かもしれません。電磁気動力学が発展する以前には、電静力学の法則は電気の法則として見做されていました。現在では、電荷間および座標系に対する相対的に静止している場合にのみ、電静力学から正確な電場を導くことが可能です。このため、マクスウェルの方程式が電磁気動力学における電静力学を覆す理由があるのでしょうか？全くありません。電静力学は電磁気動力学の中に限界として含まれており、後者の法則は時間に不変の電場の場合には直接的に前者の法則へと導かれます。物理理論に与えられる最も公正な運命は、それがより広範な理論への導入方法を自ら示すことです。その結果として、それ自身が新たな理論の中で生き残る限界となるのです。

光の伝播に関する例で見たように、一般相対性理論は重力場によって自然過程の進行に及ぼす影響を理論的に導き出すことができます。重力場がない場合での法則は既知であるときです。しかし、最も魅力的な問題の一つは、一般相対性理論が鍵を提供するものであり、それは自身の重力場を満たす法則の調査に関連しています。この点について一瞬考えてみましょう。

空間時間領域は、適切な参照系の選択により「ガリレオ的」に振る舞う（近似して）ものと知られている。すなわち、重力場が存在しない領域である。もしもこのような領域を、どのような運動でも持つ参照系 \( K' \) に参照すると、\( K' \) 相対的に重力場は空間および時間に変動するものが存在する。この重力場の性質は、選択された \( K' \) の運動によって当然異なるものとなる。一般相対性理論によれば、このような特別な種類の重力場が得られるすべての重力場に対して、重力場の一般的な法則を満たすべきである。確かにすべての重力場がこの方法で生成されるわけではなくても、重力場の一般的な法則は、そのような特別な種類の重力場から導き出せるという希望を持つことができる。この希望は最も美しい形で実現された。しかし、この目標に対する明確な視野とその実現間には重要な困難があり、それが根本的なところに存在するため、読者からそれを隠すことは勇気が必要である。私たちは空間時間連続体の概念をさらに拡張する必要がある。


\clearpage
\chapter{第XXIII節：Behaviour of Clocks and Measuring-Rods on a Rotating Body of Reference}

XXIII.

回転座標系における時計と測定棒の挙動

これまで、一般相対性理論における空間・時間データの物理学的解釈について言及してこなかった。そのため、私たちは特別な理論の相対性原理から離れて一定の怠慢さを示している。これは、特別相対性原理から学ぶように、決して無視できるほど重要でないものではない。今やこの欠点を補う時期が来ている；しかし最初に述べておくと、この問題は読者の忍耐力と抽象化能力に大きな要求をかけている。

私たちは再び非常に特別なケースから始めよう。これらは以前にもしばしば使用した。空間時間領域を考えるが、その領域において相対性体 $K$ に対して重力場が存在しないとする。そして $K$ の運動状態が適切に選ばれているとしよう。そうすると、考慮されている領域に関して $K$ はガリレオ座標系であり、特別相対性原理の結果は $K$ 相対的に成り立つ。次に、同じ領域を回転座標体 $K′$ に対して考えるとする。これは $K$ に対する回転運動が一様であるとする。具体的には、平面円盤状の $K′$ を想像し、その中心を中心に自身の平面上で一様な回転を行うとしよう。円盤上に偏心して座っている観測者は、放射方向に外向きの力を感じ、これは静止した原座標体 $K$ に対する観測者が「慣性効果（離心力）」として解釈するだろう。しかし、円盤上の観測者は自分の円盤を「静止している」と考える可以性がある。一般相対性原理に基づいて、彼はこれを行う正当な理由がある。自身および円盤上に静止している他の物体に対する作用する力は、重力場の効果と見做す。しかし、観測者は一般相対性理論を信じているため、これは彼を混乱させるものではない；彼は自分自身が体験する力を説明するために、一般的な重力法則を導出できると考える権利がある。この重力法則は星の運動だけでなく、自身が経験する力場も説明できると信じることができる。

 盤の中心では場がなくなり、中心から外側に向かって距離に比例して増加する。

観測者は円盤に対して時計と測定棒を使って実験を行い、円盤 \( K' \) における時間・空間データの正確な定義を得ることを意図している。これらの定義は彼の観察に基づいている。この作業において彼は何を感じるだろうか？

まず最初に、二つの同一構造の時計のうちの一つを円盤の中心に置き、もう一つは円盤の端に置き、これらは円盤に対して静止しているようにする。ここで私たちは、これらの時計が非回転するガリレオ座標系 \( K \) から見たとき、どちらの時計も同じ速度で進むかどうかを問う。この体から見れば、中心に置かれた時計は速度を持たない一方で、円盤の端にある時計は円盤の回転により \( K \) 相対的に動いている。第12節で得られた結果によると、後者の時計は常に中心に置かれた時計よりも遅い速さで進むことがわかる。つまり、\( K \) から観測すると遅い速度で進むことになる。明らかに、円盤の中心にある時計の隣に想像した観測者がその時計を見ている場合でも同じ効果が確認されるだろう。したがって、円盤上では、またはより一般的には任意の重力場において、時計は置かれた位置によって速さが異なる。そのため、参照体に対して静止している位置に配置された時計を使って時間の定義を合理的なものにするのは不可能である。また、このような状況で以前の同時性の定義を適用しようとすると同様の困難が生じるが、この問題についてさらに詳しく説明するつもりはない。

さらに、この段階では空間座標の定義自体にも不溶な困難が存在する。観測者が円盤の端に接して短い標準計量尺（半径と比較して短い）を適用すると、ガリレオ系から見れば、その長さは1未満となる。これは第十二節で述べたように、運動している物体は運動方向に縮小するからである。一方で、円盤の半径方向に計量尺を適用した場合、$K$ からはその長さは経験しない。そこで観測者が円盤の周囲と直径をそれぞれ計量尺で測定し、それを割ってπを求めると、通常の3.14...ではなく大きな数値が得られる。もちろん、円盤が静止している場合、この操作は正確にπを与える。これはユークリッド幾何学の命題が回転する円盤上で厳密には成り立たないことを証明し、一般に重力場の中でも大抵例外なく成り立たないことを示す。したがって直線という概念そのものにも意味が失われる。従って、特殊相対性理論で用いた方法を使って円盤に対する座標 $x, y, z$ を正確に定義することは不可能である。そして座標と事件の時間が定義されない限り、これらの法則を正確な意味を持つ形で表現することもできない。

 ここでの考察においては、ガリレオ（回転していない）系 $K$ を参照体系として用いる必要がある。これは特殊相対性理論の結果が $K$ （または $K′$ では重力場が存在する）に対してのみ有効であると仮定できるからである。

したがって、我々が以前に導き出した一般相対性理論に基づく結論は疑問視される可能性がある。実際には、この問題を正確に扱うために微細な迂回が必要となる。私は次の段落でこれを準備する。


\clearpage
\chapter{第XXIV節：Euclidean and Non-Euclidean Continuum}

ユークリッドと非ユークリッド連続体

大理石のテーブルの表面が私に広がっている。このテーブル上の任意の一点から他の任意の一点へ移動するには、一点から「隣接」の一点を経由して移動し、これを繰り返すことで実現できる。言い換えれば、点から点へと移動しながら「ジャンプ」は行わないことである。読者にはここでの「隣接」と「ジャンプ」について十分に理解していただけると思う（過度に論理的にならない限り）。この性質を表面が連続体であるという表現で表す。

次に、等しい長さの小さな棒が多数作られたと想像する。これらの棒の長さは大理石板の寸法に対して非常に小さいものとする。棒の長さが等しいとは、一つを他の任意の棒に並べて端が重ならないことを意味することである。その後で、四本のこれらの小さな棒を marble board に配置し、正方形を形成するようにする。その対角線は等しくなるようにするためには、小さな検定棒を使用する。この正方形に似た他の正方形を追加し、最初の正方形と共通の一本の棒を持つようにする。これらの各正方形について同じ手順を繰り返すことで、最終的に marble board は正方形で覆われる。その配置は、各正方形の辺が二つの正方形に属し、それぞれの角が四つの正方形に属することを特徴とする。

この作業を成功させるには多くの驚きが必要である。例えば、どの瞬間でも三つの正方形が一点で接する場合、第四の正方形の二辺は既に配置されており、残りの二辺の配置は完全に決定される。しかし、その結果、四角形の対角線が等しくなるように調整することは不可能となる。もし等しいと自然に成るなら、これは大理石板と小さな棒に対する特別な恵みであり、私はそれを感謝することができる。この構築が成功するためには多くのこのような驚きを経験しなければならない。

もし本当に順調に進んだなら、大理石の板上の点は小さな棒を「距離」として使用した場合、「ユークリッド連続体」であると言えるでしょう。一方の隅を「原点」と定めると、他の隅はこの原点に対する2つの数で特徴付けられます。私はどのように多くの棒を通らなければならないかを述べるだけで、原点から「右」へ進み、「上」に移動して、考慮中の隅に到着することができるでしょう。これらの2つの数字は、小さな棒の配置によって決定される「カartesian座標系」におけるこの隅の「カartesian座標」です。

この抽象的な実験を以下のように変更することで、実験が失敗することも認識できます。我々は小さな棒が温度上昇に比例して拡大するという仮定をします。中央部の大理石の板を加熱し、周囲は加熱しないとしましょう。この場合でも、テーブル上の任意の位置で2つの小さな棒を一致させることができます。しかし、加熱により、テーブルの中心部にある小さな棒が拡大する一方で、外周の小さな棒は変化しないため、正方形の構築は混乱してしまいます。

我々の小さな棒（単位長度として定義されたもの）に関して言えば、大理石の板は Longerユークリッド連続体」ではなくなり、カartesian座標を直接作成する立場からも脱却せざるを得ません。しかし、テーブルの温度に似た影響を受けない（あるいは全く影響を受けない）他の物があるため、大理石の板が「ユークリッド連続体」と見なすことは自然です。これは、長さの測定や比較についてより精巧な規定を設けることで達成できます。

ただし、すべての種類の棒（すなわち、どの素材でも）が温度変化に対して大理石の平板上に置かれたときに同じように行動し、我々が温度の影響を検出する唯一の方法が、上記で説明したような実験において棒の幾何学的挙動から得られるものであると仮定すれば、最も良い計画は、棒の一端を平板上の二点に一致させる距離を定義することとなる。なぜなら、それ以外では距離を定義する方法がいかに極めて恣意的であるか。デカルト座標法はこの方法を放棄し、ユークリッド幾何学を仮定しない別の方法に置き換えるべきである。

読者には注意していただきたいが、ここに描写された状況は一般相対性原理（第二十三節）によって引き起こされる状況と一致している。

 数学者たちは我々の問題を以下のように遭遇した。与えられた曲面（例えば、楕円体）がユークリッド三次元空間に属するとき、その曲面上には二次元幾何学があると同様に平面曲面上にもあるように、二つの次元の幾何学が存在する。ガウスはこの二次元の幾何学を最初から扱う方法を探し求め、三次元ユークリッド空間の事実を利用せずにその問題を解決した。棒で構築物を曲面上に作る（大理石の平板の場合と同様）と、平面幾何学に基づく法則とは異なる法則が得られることがわかるだろう。曲面は棒に対してユークリッド三次元空間の連続体ではないので、曲面上ではデカルト座標を定義することができない。

ガウスは曲面上での幾何学的関係を扱うための原理を示し、リーマンが扱った多次元非ユークリッド連続体の方法への道を開いた。したがって、数学的には既に我々が一般相対性原理から導かれる形式的な問題を解決していた。

\[ \]


\clearpage
\chapter{第XXV節：Gaussian Co-ordinates}

ガウスによれば、この解析と幾何学を組み合わせた手法は次の方法で達成できる。テーブルの表面に任意の曲線（図4参照）を描くことを想像する。これらを$u$-曲線と呼び、それぞれの曲線を数で示す。$u$ = 1, $u$ = 2, $u$ = 3 の各曲線が図において描かれている。$u$ = 1 と $u$ = 2 の間には、1と2の間に実数として存在するすべての実数に対応する無限に多くの曲線を想像する必要がある。これにより$u$-曲線の系ができ、これは「無限に濃密」な系であり、テーブルの全体表面を覆う。これらの$u$-曲線は交差しないものでなければならないし、表面の各点には一つだけ曲線が通る。したがって、 marble slab の表面のどの点にも完全に定まった$u$の値が存在する。同様に、表面に$v$-曲線を想像する。これらも$u$-曲線と同様の条件を満たし、それぞれに数が割り当てられる。また、任意の形状であってもよい。したがって、表面のどの点にも$u$の値と$v$の値が存在する。これらの二つの数をテーブルの表面の座標（ガウス座標）と呼ぶ。例えば、図における点Pはガウス座標で$u$ = 3, $v$ = 1である。表面の近接点PとP′には次のような座標が対応する：

$P$: $u, v$

$P'$: $u$ + $du, v$ + $dv$,

ここで、$du$と$dv$は非常に小さな数を表す。同様に、棒で測定した点PとP′の距離（線間隔）を非常に小さな数$ds$によって示すことができる。ガウスによれば、

\[ ds^2 = g_{11} du^2 + 2 g_{12} du\, dv + g_{22} dv^2 \]

が成り立つ。ここで、$g_{11}, g_{12}, g_{22}$は$u$と$v$に完全に定まった方法で依存する量である。これらの量$g_{11}, g_{12}, g_{22}$は棒の挙動を$u$-曲線および$v$-曲線に対して、そして表面に対して決める。表面の点が測尺に対するユークリッド連続体である場合のみ、すなわちこの場合のみ、$u$-曲線と$v$-曲線を描き、これらに数を割り当てることができるようになる。その結果、

\[ ds^2 = du^2 + dv^2 \]

となる。

以下の条件のもとでは、$u$-曲線と$v$-曲線はユークリッド幾何学における直線であり、これらは互いに垂直です。ここでのガウス座標は単なる直交座標です。ガウス座標は表面の点に対し非常に僅かな数値の差異を持つ隣接する点と関連付けるような二つの数値セットとの関連付けに過ぎません。

以上までで、これらの考慮事項は2次元の連続体に対して成り立ちます。しかし、ガウス法は3次元、4次元以上の連続体にも適用できます。例えば、4次元の連続体が利用可能と仮定すると、以下のように表現できます。連続体の各点に対し、任意に四つの数値$x_1, x_2, x_3, x_4$を関連付けることができます。これらは「座標」と呼ばれます。隣接する点は座標の隣接する値に対応します。隣接する点$P$と$P'$間に距離$ds$が関連付けられ、物理的に測定可能な明確に定義されたものである場合、以下の式が成り立ちます：

\[ ds^2 = g_{11} dx_1^2 + 2 g_{12} dx_1 dx_2 + \dots + g_{44} dx_4^2 \]

ここで、$g_{11}$ 等の量は連続体内での位置に依存して変化します。ユークリッド連続体である場合のみ、座標$x_1, \dots, x_4$ が連続体の点と関連付けることができ、以下のように表現できます：

\[ ds^2 = dx_1^2 + dx_2^2 + dx_3^2 + dx_4^2 \]

この場合、4次元連続体内で関係が存在し、これらは我々の三次元測定における関係と類似しています。

しかし、上記で与えられた$ds^2$ のガウス処理は常に可能ではありません。十分に小さな領域において考慮される連続体をユークリッド連続体として扱うことができる場合のみ可能です。例えば、テーブル上の大理石板と局所的な温度変化の場合が明らかに該当します。板の一部では温度は実質的に一定であり、したがって棒の幾何学的振る舞いはユークリッド幾何学の規則によるとほぼ正しくなります。そのため、前節で述べた四角形の構築の欠点は、その構築がテーブル表面の大半を覆うまで明らかに現れません。

ガウスは、一般の連続体に対する数学的な扱い方を発明しました。これは、「サイズ関係」（隣接点間の「距離」）を定義することから始まります。
連続体の各点には、その連続体の次元数分の数字（ガウス座標）が割り当てられます。この割り当て方は、一つだけ意味を持つようにされ、隣接する点には微小に異なる数字（ガウス座標）が割り当てられるようにします。
ガウス座標系は、カartesian座標系の論理的な一般化です。また、非ユークリッド連続体にも適用可能ですが、考慮している連続体の小さな部分がよりユークリッド系に近づくほど、その小さな部分が注意の対象となるほど、定義された「サイズ」または「距離」に関して適用可能です。


\clearpage
\chapter{第XXVI節：The Space-Time Continuum of the Special Theory of Relativity Considered as a Euclidean Continuum}

XXVI.

特殊相対性理論における空間時間連続体をミンコフスキーのユークリッド連続体として考慮する

現在、私たちはミンコフスキーのアイデアをより正確に述べることができる位置にあります。これは第XVII節でしか明確には示されていませんでした。特殊相対性理論によれば、四次元空間時間連続体の説明のために特定の座標系が優先されます。これらは「ガリレオ座標系」と呼びました。これらのシステムでは、事件を定める4つの座標 $x, y, z, t$ は、本書の最初の部分で詳細に述べたように物理的に単純に定義されています。一方のガリレオ座標系から別のガリレオ座標系への移行においては、ローレンツ変換の式が有効です。これらの式は特殊相対性理論からの推論の基礎となりますが、それ自体はすべてのガリレオ座標系に対する光の伝播法則の普遍的な有効性を表現するものに過ぎません。

ミンコフスキーは、ローレンツ変換が次の単純な条件を満たすことを発見しました。2つの近接した事件について考えます。これらのイベントの相対的位置はガリレオ座標系 $K$ に対する空間座標差 $dx, dy, dz$ と時間差 $dt$ で与えられます。一方の別のガリレオ座標系に対して、これら2つのイベントに対する対応する差分を \\textit{dx′, dy′, dz′, dt′} とします。これらの量は常に次の条件を満たします。

\[
dx^2 + dy^2 + dz^2 - c^2 dt^2 = dx'^2 + dy'^2 + dz'^2 - c^2 dt'^2.
\]

この条件からローレンツ変換の有効性が導けます。これについて次のように表現できます：四次元空間時間連続体上の2つの隣接する点に属する量

\[
ds^2 = dx^2 + dy^2 + dz^2 - c^2 dt^2,
\]

は選択されたすべての（ガリレオ）座標系に対して同じ値を持ちます。 $x, y, z$ を $x_1, x_2, x_3, x_4$ に置き換えると、次のようにもなります：

\[
ds^2 = dx_1^2 + dx_2^2 + dx_3^2 + dx_4^2.
\]

参照系の選択によらず独立している。これを「二つの事象または四次元点の間の『距離』」と呼ぶ。

したがって、実数量 $t$ ではなく虚数変数

\[
\tau = \frac{ict}{\sqrt{1 - v^2/c^2}}
\]

を時間変数として選ぶ場合、特殊相対性理論に従って空間時連続体は「ユークリッド」四次元連続体とみなすことができる。これは前節の考慮から導かれる結果である。


\clearpage
\chapter{第XXVII節：The Space-Time Continuum of the General Theory of Relativity is Not a Euclidean Continuum}

XXVII.

一般相対性理論の空間時間連続体はユークリッド的なものではない

本書の最初の部分では、物理的解釈が容易で直接的な空間時間座標を使用することができました。そして第26節によれば、この座標系は四次元直交座標として扱えると見なすことができます。これは光速の恒常性の法則に基づいていました。しかし、第21節では一般相対性理論がこの法則を保持できないことを示しました。逆に、我々はその理論によれば、重力場がある場合、光速は座標に依存するものであるという結果を得ました。第23節の特定の例と関連して、重力場の存在が座標と時間の定義を無効にすることが見つかり、これは特殊相対性理論の目的につながりました。

これらの考慮結果から、一般相対性原理に基づいて空間時間連続体はユークリッド的なものではなく、そこでは温度の局所的変動を持つ大理石板のような一般の場合があると結論付けることができます。それは第23節で例として紹介した二次元連続体と同じです。 marble slabに例えたように、等しい棒から直交座標系を構築することは不可能でしたように、ここでも測定棒や時計から一様な配置のシステム（参照物体）を作り上げることは不可能です。それは位置と時間を直接的に示すものでなければなりません。それは第23節で我々が直面した困難の本質でした。

しかし、第25節と第26節での考慮結果はこの困難を克服する道を開きます。四次元空間時間連続体をガウス座標に任意に参照します。連続体（事件）の各点に対して4つの数値 $x_1, x_2, x_3, x_4$（座標）を割り当てます。これらの数値は全く直接的な物理的意味を持たないものであり、ただ連続体の点を明確かつ任意の方法で番号付けするための手段です。この配置は、 $x_{1}$, $x_{2}$, $x_{3}$ を「空間」座標とし、 $x_{4}$ を「時間」座標と見なす必要があるものであるとは限らないという点に注意してください。

読者には、このような世界の説明が不十分であると考えられるかもしれません。イベントに特定の座標 $x_1, x_2, x_3, x_4$ を割り当てるとは何を意味するのでしょうか？これらの座標自身は意味を持たないのですから。しかし、慎重な考慮によれば、この懸念は根拠がないことが明らかになります。

例えば、任意の運動を持つ物質点を考えましょう。この点が時間にわたって存在しないなら、その空間・時間での描写は一組の座標値 $x_1, x_2, x_3, x_4$ で行われます。したがって、その永久的な存在は無数のこのような座標値のシステムによって特徴付けられるべきであり、これらの座標値は連続性をもたらすほど密接に並んでいます。物質点に対応するものは、四次元連続体における一維線です。同様に、我々の連続体におけるこのような線は動く多くの点に対応します。これらの点に関する物理的な存在を主張できる唯一の声明は、実際にはその遭遇に関する声明です。数学的な扱いにおいて、このような遭遇は、所問の点の運動を表す二つの線が共通の座標値 $x_1, x_2, x_3, x_4$ をもつという事実によって表現されます。熟考した読者はおそらく、物理的な声明において我々に遭遇する時間空間性の唯一の実際的証拠はこのような遭遇であることに同意することでしょう。

物質点が参照体に対する運動を説明する際に述べたことは、その点が特定の参照体上の点と遭遇したという事実だけでした。時計との遭遇を観察することで、対応する時間値も決定できます。時計の針が特定の刻度点と遭遇する際の遭遇も観察します。測量棒を使用して空間測定を行う場合も、少々の考慮から明らかになります。

以下の声明は一般に成り立つ：物理的な説明は、それぞれの文が二つの事件AとBの空間時間の一致を参照する数多くの声明に分解される。ガウス座標においては、これらの各々の声明はその四次元座標 \(x_1, x_2, x_3, x_4\) の一致によって表現される。したがって実際には、ガウス座標を用いて時空連続体を説明することは、参照系を用いた説明から脱却し、その欠点から解放され、連続体のユークリッド的な性格にとらわれることなく完全に置き換えることができる。


\clearpage
\chapter{第XXVIII節：Exact Formulation of the General Principle of Relativity}

XXVIII.

一般相対性原理の正確な公式化

現在、第十八節で与えられた暫定的な一般相対性原理の公式を正確な公式に置き換えることができます。そこで用いられた形、「すべての参照系 $K, K′$, などは自然現象（自然法則の表現）の説明においてどのような運動状態であれ同等である」というものは、空間時間記述において剛体参照系を使用することは一般には不可能であるため、維持することができません。ガウス座標系が参照系として用いられるべきです。以下に述べる声明は、一般相対性原理の基本的なアイデアに対応します。「すべてのガウス座標系は自然法則の表現において本質的に同等である」。

さらにこの一般相対性原理を別の形で表すこともできます。これは特殊相対性原理の自然な拡張形式よりもより明確に理解できる形です。特殊相対性原理によれば、一般の自然法則を表現する方程式は、ローレンツ変換を使用して（ガリレオ参照系 $K$ の）空間時間変数 $x, y, z, t$ を新しい参照系 $K′$ の空間時間変数 \\textit{x′, y′, z′, t′} に置き換えるとき、同じ形式の方程式になります。一方、一般相対性原理によれば、ガウス変数 $x_1, x_2, x_3, x_4$ の任意の置換によって、方程式は同じ形式のものに変化する必要があります；すべての変換（ローレンツ変換だけではなく）は一つのガウス座標系から別のガウス座標系への移行に対応します。

もし三次元的な視点を維持したいのであれば、一般相対性理論が経験している発展は次のように特有な概念で説明できます：特別相対性理論はガリレオ空間に適用されます、すなわち、重力場の存在しない領域に適用されます。このような状況では、ガリレオ参照体が参照系として機能し、つまり、「孤立」した物質点に対する等速直線運動のガリレオ法則がその上に成り立つように選ばれた剛体です。

ある考慮事項から、これらのガリレオ空間も非ガリレオ的な参照体に対して適用されるべきであると考えられます。それらの参照体に関して特別な種類の重力場が存在します（第XX節と第XXIII節を参照）。

重力場内では、剛体という概念自体はユークリッド性を持たないものとして認識されます；したがって、仮想的な剛体参照は一般相対性理論において役立つものではありません。時計の運動も重力場によって影響を受け、これにより直接時計を使用して定義された物理的時間の説明は特別な相対性理論におけるそれほど信頼性が高くありません。

このため、柔軟な参照体が使用されます。これらは全体として任意の方法で動くだけでなく、運動中に形状を任意に変更することも可能です。どのような不規則な運動法にも従う時計を使用して時間の定義を行います。これらの時計は非柔軟な参照体上の各点に固定されています。これらの時計には一つだけ条件があります：空間上隣接する時計で観測される「読み取り」が、無限小の差異を持つことです。この非柔軟な参照体は、「モルシス」と呼ぶことが適切かもしれません。これは主に任意に選択されたガウス四次元座標系と同等です。「モルシス」がガウス座標系よりも一定の理解性をもたらすのは、個々の存在としての形式的な保持が実質的に正当化されていないからです。

空間座標に対し時間座標を用いる。 mollusca 上の任意一点は空間点とみなされ、その相対的に静止している任意の物質点も同様に静止とされる。mollusca を参照体として考える限りにおいてである。一般相対性原理は、これらの mollusca が自然法則を表す際に等しく使用できる権利を持ち、また成功する可能性があることを要求する。自然法則自体は、参照体の選択に依存してはならない。

一般相対性原理の偉大な力は、上記で見たようにそれが自然法則に対して課す広範かつ厳格な制限から来ている。


\clearpage
\chapter{第XXIX節：The Solution of the Problem of Gravitation on the Basis of the General Principle of Relativity}

XXIX.

重力問題の解決に関する一般相対性原理に基づく方法について、読者はこれまでの議論を踏まえれば、さらなる困難はないと理解できるでしょう。

私たちは、ガリレオ系$K$において重力場が存在しない領域について考える始めます。特別相対性理論から、測定棒と時計が$K$に対する振る舞いは既に知られていますし、“孤立”な物質点の振る舞いも同様です；これらは直線運動を続けて均一に移動します。

次に、この領域をランダムなガウス座標系や“モルシュー”$K′$としての参照体に参照し、重力場$G$（特定の種類）が存在します。測定棒と時計、そして自由に移動する物質点の$K′$に対する振る舞いは単なる数学的変換によって学びます。これらの振る舞いを重力場$G$の影響によるものとして解釈します。

その後、仮説を導入します：重力場が存在する場合でも、測定棒、時計および自由に移動する物質点に対する重力場の影響は、ガリレオ特別の場合から直接導出されない場合でも、座標変換だけで同じ法則に従うと仮定します。

次に、ガリレオ特別の場合から単なる座標変換によって導き出した重力場$G$の空間時間的な振る舞いを調査します。この振る舞いは、説明において使用される参照体（モルシュー）を選択方法に関係なく常に有効な法則として形式化されます。

この法則はまだ一般の重力場の法則ではありません；現在考慮している重力場は特別な種類です。一般の重力場の法則を発見するためには、上記の法則を一般的化することが必要で、これは以下の要請に従って無恣意的に得られます：

(a) 必ずしも要求される一般化は、相対性原理の一般の公理にもとづいています。

($b$) 考慮される領域に存在する任意の物質は、その慣性質量のみであり、従って第XV節においてはそのエネルギーのみが場を興奮させる効果にとって重要である。

($c$) 重力場と物質とはエネルギー（および動量）保存法則を満足するべきである。

一般相対性原理により、重力場の影響を、重力場が存在しない場合に既知の法則に従って進行するすべての過程に対し決定できる。これは測長具、時計、自由に動く物質点に対する方法と同じ原則に基づいて進められる。

このように導き出された相対性原理から得られた重力理論は、美しさだけでなく、第XXI節で明らかになったクラシック力学の欠陥を解消し、慣性質量と重力質量の等価性という実験法則を解釈する点でも優れている。しかし、古典力学では無力な観測結果を説明できていない。

理論の適用を、重力場が弱く、すべての質量が座標系に対する速度が光速と比較して小さい場合に限定すれば、最初の近似としてニュートンの理論を得ることができる。このようにニュートンの理論は特別な仮定なしで得られることになるが、ニュートンは相互作用する物質点間の吸引力が距離の二乗に反比例することを仮説として導入しなければならなかった。計算精度を高めれば、ニュートンの理論からの逸脱が現れるが、これらの大半は観測によって検証されないほど小さいためである。

ここでは、これらの異常の一つについて注意を喚起しなければなりません。ニュートンの理論によれば、惑星は太陽を中心に楕円軌道上を動きますが、固定星の運動や他の惑星の影響を無視すれば、その楕円軌道は固定星系に恒久的に位置を保つべきです。したがって、これらの二つの影響を考慮に入れて観測された惑星の動きを修正し、ニュートンの理論が厳密に正確であると仮定すると、惑星の軌道は固定星系に対する静止楕円となるはずです。この推論は非常に精密に検証できるものであり、現在までの観測精度で得られる精度において、すべての惑星について確認されていますが、一つだけ例外があります。その例外は太陽に最も近い惑星である水星です。レヴィエール時代から知られているように、水星の軌道に対応する楕円を上記の影響を修正した後でも、固定星系に対する静止状態ではなく、非常にゆっくりと軌道平面内で軌道運動と同じ方向に回転していることが判明しました。この回転運動の値は1世紀あたり43秒分の弧であり、その正確さは数秒以内まで保証されています。この効果を古典力学で説明するには、非常に低確率であるとされる仮定が必要となりますが、それらはこの目的のために作られたものでした。

一般相対性理論に基づいて見れば、惑星の軌道が太陽周りに回る楕円も必然的に上記のように回転することを見つけることができます。現在までの観測精度では微細な変化しか検出できないため、すべての惑星についてこの回転は確認されていませんが、水星の場合には1世紀あたり43秒分の弧となる結果があり、これは観測と厳密に一致しています。

除くこの1つ以外、理論から観測によって検証できる2つの推論のみが可能でした。すなわち、太陽の重力場による光線の曲率と、大規模な星から到達する光のスペクトル線の移動は、類似の方法で地上に生成された対応する線と比較して発生します。これらの理論からの2つの推論は両方とも確認されました。

 1919年にエディントンらによって初めて観測されました。（付録IIIを参照）

 1924年にアダムスによって確立されました。（p.132を参照）


\clearpage
\part{宇宙全体に関する考察（Considerations on the Universe as a Whole）}



\clearpage
\chapter{第XXX節：Cosmological Difficulties of Newton's Theory}

XXX.

ニュートン理論の宇宙論的困難

セクションXXIで議論された難点以外に、クラシック天体力学にはもう一つ根本的な困難があります。この困難は、私の知識では最初に詳細に議論した天文学者セリジェによって指摘されました。全宇宙をどのように考えればよいかという問いを考えたとき、まず思い浮かぶ答えはおそらくこれでしょう：空間（および時間）に関して言えば、宇宙は無限大です。星はどこでも存在しており、密度は詳細によりますが平均的にはどの場所でも同じです。言い換えれば：どんなに遠くまで空間を旅しても、類似の種類と密度を持つ薄い恒星群を見つけることでしょう。

この見解はニュートンの理論とは調和しません。その理論によれば、宇宙には星の密度が最大となる中心点があり、そこから外側に向かって星群の密度が減少し、最終的には遠方で無限に広がる空っぽの領域に取って代わられるべきです。恒星の宇宙は空間という無限の海洋の中の有限な島地帯であるべきです。

証明—ニュートンの理論によれば、「力線」の数は無限遠方から質量mに終焉するものはその質量mに比例します。平均して、宇宙全体で密度ρ0が一定であれば、体積Vの球体内には平均的な質量ρ0\\textit{Vが含まれます。したがって、球面Fを通じて内側へと通過する力線の数はρ0}Vに比例します。球面の表面積1平方単位あたりに入り込む力線の数はρ0\\textit{V/Fまたはρ0}Rに比例します。したがって、球の半径Rが増加すると、球面での場強度は無限大になることになりますが、これは不可能です。

この概念自体は非常に満足いくものではありません。さらに、その結果として恒星から放出される光や星系内の個々の恒星が無限空間に永遠に流出し、他の自然物体と再び相互作用することなく消えてしまうという結論を導きます。有限な物質宇宙は徐々にかつ系統的に貧困化されることになります。

この難局から脱出するため、セリジェはニュートンの法則を修正することを提案しました。彼は遠距離において質量間の引力が逆平方律より速く減少すると仮定しました。これにより、物質の平均密度がどこでも一定であり、無限大まで続くことになりながら、無限に大きな重力場は生じません。このようにして、物質宇宙が中心的な性質を持つべきという不快な概念から解放されます。

もちろん、我々は基本的な困難を克服するためにニュートンの法則を修正し、複雑化させることでその代償を払います。しかし、その修正と複雑化には実験的もしくは理論的な基礎がありません。無数の法則が同じ目的を達成することができ、なぜ一つの法則が他の法則よりも優れているのかを説明する理由は存在しません；これらのどの法則もニュートンの法則と同じくより一般的な理論的原理に基づいていません。


\clearpage
\chapter{第XXXI節：The Possibility of a "Finite" and Yet "Unbounded" Universe}

XXXI.

有限でかつ「無限」な宇宙の可能性

宇宙の構造に関する推測は、全く異なる方向に進展しています。非ユークリッド幾何学の発展により、空間の無限性を疑問視できることが認識されました。これは、思考の法則や経験と矛盾することはありません（リーマン、ヘルムホルツ）。これらの問題はヘルムホルツとポインカールによって詳細かつ卓越した説明がなされており、私はここでは触れることに留めます。

まず第一に、二次元空間における存在を想像します。平らな生物とその平らな道具、特に平らで剛性のある測定棒は、平面内で自由に動くことができます。彼らにはこの平面以外の何ものも存在しません：自分たちが経験する現象や、平らな「物」が起こすことは、彼らの平面における全ての実在です。特に、平面ユークリッド幾何学の構築は棒を使って行うことができます（例えば、第XXIV節で述べた格子構築）。これらの生物の宇宙は二次元でありながら、私たちの宇宙と同様に無限に広がっています。彼らの宇宙には棒で組み立てた無数の同じ大きさの正方形があり、すなわちその体積（面積）は無限です。このような生物が「平面」であると言った場合、それは意味のある主張であり、彼らは棒を使って平面ユークリッド幾何学の構築を行うことができるからです。この際、個々の棒は位置に関係なく常に同じ距離を表します。

次に、平面ではなく球面上の二次元存在について考えましょう。この平らな生物は測量棒その他の物体がこの表面にぴったりと収まり、脱出することができません。彼らの観察可能な宇宙は球面全体に限定されています。

これらの生物は、自らの宇宙の幾何学を平面幾何学として捉え、また測量棒を「距離」の実現として認識できるでしょうか？それはできません。なぜなら、彼らが直線を実現しようとすると、曲線を得ることになります。私たる三次元存在はこれを大円と呼びます、すなわち、定数の有限長を持つ自閉的線で、測量棒を使用して測定することができます。

同様に、この宇宙には測量棒を使って構築した正方形との面積を比較できる有限の面積があります。この考慮から得られる魅力は、これらの生物が住む宇宙は有限でありながら限界がないという事実を認識することです。

しかし、球面上の存在は世界一周する必要なく、自らの「世界」のどの部分でも、測量棒を使用しない限り、彼らがEuclid的宇宙に生きているわけではないことを確認できます。点から出発し、三次元空間で判断した直線（円弧）を等長に引きます。これらの線の自由な端点を結ぶ線を「円」と呼びます。

平面において、円の周囲長と直径の比（両方の長さは同一の測量棒を使用して測定）はEuclid的平面上の幾何学によれば一定の値πに等しく、この値は円の直径に関係なく恒常です。一方で、彼らの球面上ではこの比は

\[ \frac{C}{d} = 2\pi \]

となります。

$i.e.$ 半径が「世界球」の半径 \( R \) と比較して大きいほど、πより小さな値であり、その差はより大きな円の半径でより顕著となる。この関係を利用して、球体の生命体は、彼らの測定に利用できる「世界球」の一部が非常に小さい場合でも、彼らの宇宙（「世界」）の半径を決定することができる。しかし、この部分が極めて小さい場合、彼らは自己が球面の「世界」にいることを証明できなくなる可能性がある。なぜなら、球面上の小さな部分は同じ大きさの平面部分からわずかにしか異なるためである。

したがって、太陽系が球形宇宙の一部に占める割合が極めて小さい場合、球面生命体は自己が有限の宇宙にいるのか無限の宇宙にいるのかを決定する手段を持たない。彼らがアクセスできる「宇宙の部分」はどちらの場合も実質的に平面的またはユークリッド的なためである。

この議論から直接導かれるのは、球体生命体にとって円の周長は半径とともに増加し、最終的には「世界球」の周長に達すると、その後徐々に減少してゼロへと収束することである。この過程において、円の面積はますます増大し、最後には「世界球」全体の面積と等しくなる。

読者各位が疑問を持たれるかもしれませんが、なぜ私たちは生命体を球面上に置いたのか他に閉曲面上に置かなかったのか。この選択は、すべての閉曲面の中で唯一、その上にある全ての点が同等であるという性質を持つ球面を選んだことによる正当化があるからです。半径 \( r \) に対する円の周長 \( c \) の比は \( r \) に依存しますが、与えられた \( r \) の値に対しては「世界球」上の全ての点で同じであり、言い換えれば、「世界球」は「恒定曲率面」であると言えるのです。

この二次元球宇宙には、リーマンによって発見された三次元球空間という三次元の類似物がある。その点はすべて等価である。これは有限な体積を持つが、「半径」（2π²R³）で定義される。

球空間を想像することは可能か？空間を想像することとは、私たちの「空間」経験の要約を思い描くこと、すなわち、「剛体」の移動を通じて得られる経験を思い描くことである。この意味で、球空間を想像することはできる。

点からすべて方向に線を引いたり、測定棒を使ってそれぞれの距離 $r$ を刻み、その自由端点が球面に存在することを考える。この球面の面積（$F$）は、測定棒で組み立てた正方形を使って測ることができる。宇宙がユークリッド空間であるなら、$F$ = 4π\\textit{r²；球空間であるなら、$F$ は常に 4π}r² より小さい。$r$ の値が増加するにつれて、$F$ はゼロから最大値へと増大し、「世界半径」によって決定されるが、さらに $r$ の値が増加すると、面積は徐々にゼロに収束する。最初の段階では、始点から放射した直線はますます遠ざかるが、後には接近し、最終的には「反対点」に戻り合って始点と重なる。「球空間全体を遍歴した」という状況である。このように考えると、三次元球空間は二次元球面に非常に類似していることが容易に理解できる。それは有限（すなわち有限の体積を持つ）で、境界がない。

なお、別の種類の曲線空間も存在する：「楕円空間」である。これは、「反対点」が同一（互いに区別できない）とみなされる曲線空間である。したがって、楕円宇宙は中心対称的な曲線宇宙として部分的に考えることができる。

それから述べたことから、境界のない閉じられた空間は考えられる。これらの空間の中では、球面空間（そして楕円形空間）が最も単純であり、その上にあるすべての点は等価である。この議論により、天文学者や物理学者に対して興味深い問題が生じる。それは我々が住む宇宙は無限であるか、あるいは球面宇宙のように有限であるか否かという問題である。私たちの経験はまだ十分でないため、この質問に答えることは難しい。しかし、一般相対性理論はこの質問に対する答えを一定程度確実なものにすることができ、その点で述べられた困難が解決される。


\clearpage
\chapter{第XXXII節：The Structure of Space According to the General Theory of Relativity}

空間の構造に関する一般相対性理論による考察

一般相対性理論によれば、空間の幾何学的性質は独立したものではなく、物質によって決定される。したがって、宇宙の幾何学的な構造について結論を導くためには、物質の状態という何か既知のものに基づいて議論を行う必要があります。経験から、適切に選んだ座標系においては、星の速度は光の伝播速度と比較して小さいことがわかります。したがって、物質を静止として扱う粗略な近似で宇宙全体の性質について結論を導くことができます。

我々はすでに前節での議論から、測量棒や時計の行動が重力場によって影響を受け、すなわち物質の分布によって影響を受けていることを知っています。これだけでユークリッド幾何学の正確な有効性が宇宙において排除されることを示すのに十分です。しかし、我々の太陽規模の質量に至るまで周囲空間の計量が極めて小さな程度にしか影響されないという計算結果から、我々の宇宙はユークリッド的なものと僅かに異なる可能性があると考えることができます。そして、このような考え方は、太陽規模の質量による影響が極めて小さいことを示す計算結果からより確実です。

我々は測量棒や時計の行動が重力場によって影響を受けていることから、幾何学的に我々の宇宙は個々の部分で不規則に曲がっている表面のように振る舞うかもしれませんが、どこでも平面から大きく離れるわけではありません。湖面のような波打った表面のようなものです。このような宇宙を擬ユークリッド宇宙と呼びましょう。その空間に関して無限大であることが示されます。しかし計算結果によれば、擬ユークリッド宇宙において物質の平均密度は必ずしもゼロではないことを示す必要があります。したがって、このような宇宙はどこでも物質で満たされていない可能性があり、我々が第30節で描いた不十分な情景を提示することになります。

以下の原著英文を忠実に日本語へ翻訳しました：

もし宇宙中の物質の平均密度がゼロとは異なる値である場合、無論その差がいかに小さくても、宇宙は準ユークリッド的でないことが必要となります。逆に計算結果は、物質が均一に分布している場合、宇宙は必ず球形（または楕円形）になることを示しています。実際のところ、物質の詳細な分布は均一ではなく、宇宙は個々の部分で球形から逸脱し、すなわち宇宙は準球形となります。しかし、それは必然的に有限です。実際には、宇宙の空間的広がりとその中の平均密度との間に関係式が提供されています。

宇宙の半径 \( R \) に対する方程式は以下のようになります：

\[
\frac{2}{k} = R
\]

この式において、\( k \) はニュートンの引力定数と関連する定数であり、平均物質密度を \( \rho \) とすると、C.G.S. 系統を使用した場合、2/k = 1.08 x $10^{27}$ となります。

付録


\clearpage
\part{付録（Appendices）}



\clearpage
\chapter{付録 I：Simple Derivation of the Lorentz Transformation}

付録 I

ローレンツ変換の単純な導出（第十一節の補遺）

図2で示した座標系の相対的な向きについて、両方のシステムの $x$-軸は常に一致します。現在の場合では、問題を部分に分割し、まず $x$-軸上に局在するイベントのみを考えることができます。このような任意のイベントは、座標系 $K$ に関して $x$ の横座標と時間 $t$、座標系 $K′$ に関して $x′$ の横座標と時間 $t′$ で表されます。 $x$ および $t$ が与えられたときの $x′$ および $t′$ を求めます。

$x$ 軸の正方向に進行する光信 hiệuは、方程式

$x$ = $ct$

または

$x$ – $ct$ = 0 . . . . . (1)

によって伝送されます。同じ光信 hiệuが座標系 $K′$ 相対的に速度 $c$ で伝送されるため、座標系 $K′$ の伝播は類似の式

$x′$ – $ct′$ = 0 . . . . . (2)

によって表されます。方程式（1）を満たす空間時間点（イベント）もまた方程式（2）を満たします。これは、一般に次の関係が成り立つときであることが明らかで、

($x′$ – $ct′$) = λ($x$ – $ct$) . . . (3)

ここで λ は定数を表します。なぜなら、(3) から ($x$ – $ct$) の消失は ($x′$ – $ct′$) の消失を意味するからです。

負の $x$-軸方向に伝送される光線について類似の考慮を適用すると、条件

($x′$ + $ct′$) = ($x + ct$) . . . (4)

を得ます。方程式（3）と（4）を加算または減算し、便利のために定数 λ および μ の代わりに定数 $a$ および $b$ を導入すると、



\[ x = \frac{bc}{a} t \]



の式が得られます。したがって、問題の解は定数 $a$ および $b$ が分かれば得られます。これらの値は以下の議論から得られます。

座標系 $K′$ の原点に関して常に $x′$ = 0 であり、従って方程式（5）の最初の式より



\[ v = \frac{bc}{a} \]



となります。 $K′$ の原点が $K$ 相対的に速度 $v$ で動いていると仮定すると、



\[ a^2 - b^2 = 1 \]



となります。同じ値 $v$ は、座標系 $K′$ の別の点が $K$ 相対的に持つ速度、または $K$ 上の点が $K′$ 相対的に持つ（負の $x$-軸方向に向かう）速度から方程式（5）を用いても得られます。つまり、 $v$ を二つの座標系の相対的速度と定義できます。

さらに、相対性原理は、Kから見れば、静止基準系$K′$に沿った単位測定棒の長さが、$K′$から見れば静止基準系$K$に沿った単位測定棒の長さと全く同じであることを教えています。$x′$-軸の点が$k$からどのように見えるかを見るためには、$k$からの$k′$の「シャッットショット」を取るだけで十分です。これは、方程式(5)の最初の式から特定の値（時間t）を挿入することを意味し、たとえば $t = 0$。

この値の$t$から方程式(5)の最初の式を使用して

\[ x' = ax \]

を得ます。したがって、$k′$系で距離Δ$x′$ = 1に分かれた2点は、瞬間的な写真では距離

\[
\Delta x = a \Delta x'
\]

に分かれます。

一方、$k′$から（t′ = 0）シャッットショットを撮影し、方程式(5)から$t$を消去して式(6)の表現を使用すると、

\[
x^2 - c^2 t^2 = x'^2 - c^2 t'^2
\]

を得ます。

この結果から、$k$系における$x$-軸上の距離1に分かれた2点は、シャッットショットでは距離

\[
\Delta x = \frac{\Delta x'}{\gamma}
\]

に表示されることを結論付けることができます。ここでγはローレンツ因子です。

したがって、両方のシャッットショットは同じであるため、(7)式におけるΔ$x$は(7$a$)式におけるΔ$x′$と等しくなければなりませんから、

\[
\Delta x = \frac{\Delta x'}{\gamma}
\]

を得ます。式(6)と(7$b$)は定数aとbを決定します。これらの定数の値を(5)に代入すると、第十一節で与えられた方程式の最初と第四の式が得られます。

したがって、$x$-軸上の事象に対するローレンツ変換を得ました。これは条件

\[
x'^2 - c^2 t'^2 = x^2 - c^2 t^2
\]

を満たします。

この結果の拡張は、$x$-軸外の事象も含めることで得られます。これには方程式(8)を保持し、

\[
\begin{aligned}
x' &= \gamma (x - vt), \\
t' &= \gamma \left(t - \frac{vx}{c^2}\right)
\end{aligned}
\]

という関係式を追加します。これにより、任意の方向を持つ光線の真空中的な速度の恒常性の公理が、基準系$K$と基準系$K′$に対して満たされます。これは次のように示すことができます。

$k$の原点から$t = 0$の時刻に光信号を送るとします。その信号は方程式

\[
x' = \gamma (x - vt)
\]

または、これを2乗すると、

\[
x^2 + y^2 + z^2 - c^2 t^2 = 0
\]

という方程式で伝播します。

以下の英文を格調高く明晰な日本語へ完全直訳しました：

光の伝播法則と相対性原理との一致により、問題の信号の伝搬は参照枠 $K′$ から見れば、

$r′$ = $ct′$

または、

$x′$2 + $y′$2 + $z′$2 – $c_{2}$$t′$2 = 0 . . . . . . (10$a$)

に従って行われるべきである。

方程式（10$a$）が方程式（10）の結果であるためには、

$x′$2 + $y′$2 + $z′$2 – $c_{2}$$t′$2 = σ ($x_{2}$ + $y_{2}$ + $z_{2}$ – $c_{2}$$t_{2}$)
(11)

が必要である。

方程式（8$a$）は $x$-軸上の点に対して成り立つため、σ = 1 となる。Lorentz変換が σ = 1 の場合に実際的に方程式（11）を満たすことは容易に確認できる；これは方程式（8$a$）と（9）の結果であるから、方程式（8）と（9）の結果でもある。したがって、Lorentz変換は導き出された。

参照枠 $K′$ の軸を $K$ の軸に平行にするか否か、または $K′$ が $K$ 相対的にどの方向に移動しているかに関わらず、Lorentz変換（8）と（9）は依然として一般化される必要がある。明らかに、参照枠 $K′$ の軸を $K$ の軸と同じ空間的方向に選ぶことは重要ではないし、また $K′$ が $K$ 相対的に移動する速度が $x$-軸方向であることも重要ではない。単純な考慮から、この一般化された意味でのLorentz変換は、特別な意味でのLorentz変換と純粋に空間的な変換の2種類の変換から構築できることが示される。

数学的には、一般化されたLorentz変換を次のように特徴付けることができる：

それは \\textit{x′, y′, z′, t′} を $x, y, z, t$ の線形同次関数として表現し、

$x′$2 + $y′$2 + $z′$2 – $c_{2}$$t′$2 = $x_{2}$ + $y_{2}$ + $z_{2}$ – $c_{2}$$t_{2}$
(11$a$)

が恒等的に成り立つようにする。すなわち：左辺に $x, y, z, t$ の表現を代入した場合、（11$a$）の左辺は右辺と一致することである。


\clearpage
\chapter{付録 II：Minkowski's Four-Dimensional Space ("World")}

Minkowskiの4次元空間（「世界」）は、ローレンツ変換をより単純に表現するため、時間変数として虚数

\[
\tau = i t
\]

を導入することができます。これに従って、

\[
x_1' = \frac{x_1 - v x_4}{\sqrt{1 - v^2}}, \quad x_2' = x_2, \quad x_3' = x_3, \quad x_4' = \frac{x_4 - v x_1}{\sqrt{1 - v^2}}
\]

を置き、同様に付加された系 \(K'\) についても同じように置き換えると、変換が恒等的に満たす条件は次のように表されます：

\[
x_1'^2 + x_2'^2 + x_3'^2 + x_4'^2 = x_1^2 + x_2^2 + x_3^2 + x_4^2 \quad (12)
\]

つまり、この「座標」の選択により、(11$a$) [Appendix II の末尾を参照] はこの方程式に変換されます。

(12)から、虚数時間座標 \(x_4\) が変換条件において空間座標 \(x_1, x_2, x_3\) と同じように入るという事実を見ることができます。これが相対性理論によれば、「時間」 \(x_4\) が自然法則の中で空間座標 \(x_1, x_2, x_3\) と同じ形で現れる理由です。

「座標」 \(x_1, x_2, x_3, x_4\) を用いて記述される4次元連続体は、Minkowskiによって「世界」と呼ばれ、点イベントを「世界点」と呼んでいました。三次元空間における物理現象は、「世界」の四次元「存在」に似た形でなります。

この四次元「世界」は、ユークリッド解析幾何学の三次元「空間」に密接な類似性があります。新しい直交座標系 \((x'_1, x'_2, x'_3)\) を導入して、同じ原点を持つ場合、\(x'_1, x'_2, x'_3\) は \(x_1, x_2, x_3\) の線形同次関数であり、

\[
x_1'^2 + x_2'^2 + x_3'^2 = x_1^2 + x_2^2 + x_3^2
\]

を恒等的に満たします。

(12)とこの類似性は完全です。Minkowskiの「世界」を形式的に4次元ユークリッド空間（虚数時間座標を持つ）として捉えることができます；ローレンツ変換は四次元「世界」における座標系の「回転」に対応します。


\clearpage
\chapter{付録 III：The Experimental Confirmation of the General Theory of Relativity}

付録 III

一般相対性理論の実験的確認

系統的な理論的な観点から見れば、実験科学の進展過程は継続的な帰納プロセスであると考えることができます。理論は観察結果を大量に集め、その形式で経験法則として表現され、比較によって一般化された法則が得られます。このように考えると、科学の発展は分類された目録を作成するような実用的な作業に似ています。これは、純粋に経験主義的であると言えます。

しかし、この観点だけでは実際の過程を完全には捉えていません。なぜなら、正確な科学における直感と論理的推論の重要な役割が省かれているからです。科学が初期段階を脱した後は、理論的な進歩は単なる整理作業ではなくなります。経験的事実に基づいて、研究者は一般的に少数の基本的な仮定（公理）から論理的に構築された思想体系を開発します。このような思想体系を「理論」と呼びます。その理論が存在する正当性は、多くの単独の観察結果と整合性を持つことにあるし、それが理論の「真実」を表しています。

同じ複雑な経験的事実に対応しても、いくつかの異なる理論があり得ます。これらの理論は相当程度異なりますが、検証可能な推論から見ると、その一致が完全である場合があります。両理論で異なる推論を見つけるのは困難になることがあります。生物学の分野では、生存競争における選択による種類の発展というダーウィンの理論と、獲得性状態の遺伝的伝達に基づく発展の理論など、一般的に興味深い例が用いられます。

Weinstein機械学の二つの理論から導き出される結論間の広範な一致のもう一つの例があります。これはニュートン力学と一般相対性理論との間で見られるものです。この一致は極めて深く、現在までに我々が発見した一般相対性理論からのわずかな数少ない導出の中でさえも、事前相対性物理学にも同様の結論を導き出すものが存在しないという状況があります。これは二つの理論の基礎的な仮定が非常に異なるにもかかわらずです。以下では、これらの重要な導出について再度考慮し、それらに関連するこれまで得られた実証的証拠についても議論します。

(a) 水星の近日点運動

ニュートン力学とニュートンの引力法則によれば、太陽を中心に回転している惑星は、その軌道が楕円形であると予測されます。正確には、太陽と惑星の共通の重心を中心として楕円を描くことになります。このような系では、太陽（または共通の重心）は軌道楕円の一つの焦点に位置し、惑星年を通じて太陽から惑星までの距離が最小から最大へと増大し、その後再び最小へと減少します。ニュートンの法則ではなく、異なる引力法則を計算に入れる場合でも、その運動は依然として、太陽から惑星までの距離が周期的に変動するような形で行われることが見出されます。しかし、この新しい法則では、一定期間（近日点から次の近日点まで）に太陽と惑星を結ぶ線が描く角度は360°から異なる値になります。その結果、軌道の線は閉じたものではなく、時間とともに太陽からの最遠点と最近点の間の円形領域を埋めていくことになります。

一般相対性理論においても、当然ながらニュートンの理論とは異なりますが、惑星の軌道運動からわずかな変動が生じると予測されます。そしてその結果として、太陽と惑星を結ぶ半径が描く角度は一つの完全な回転分を超える量により、以下のように与えられます。

\[
\Delta \phi = \frac{6\pi G^2 M^3}{c^4 a^3 (1-e^2)} t_p
\]

ここで \(G\) は万有引力定数、\(M\) は太陽の質量、\(a\) は軌道半長軸、\(e\) は偏心率、そして \(t_p\) は近日点周期です。

一方、一回の完全な回転は物理学で通常採用される絶対角測定において2πに対応し、上記の式は、一つの近日点から次の近日点までの間隔における太陽惑星距離がこの角度をどれだけ超えるかを表しています。）本式において、aは楕円の主軸半径、eはその偏心率、cは光速、Tは惑星の回転周期を表します。当該結果はさらに以下のように述べることができます：一般相対性理論によれば、楕円の主軸は太陽を中心に同じ方向で回転し、惑星の公転運動と同じ向きです。理論上は、この回転は水銀の場合に1世紀あたり43秒分の弧度となるべきであり、他の太陽系惑星についてはその大きさが非常に小さく、検出される可能性が低いと要求されます。

特に、次の惑星である金星の軌道がほぼ完全な円であるため、近日点を正確に特定することが難しいことからです。

実際には、天文学者はニュートンの理論だけでは現在の観測精度に対応するほど精密な Mercury の運動を計算することはできないと発見しました。残りの惑星による惑星への干渉効果を考慮に入れた後でも、Mercury の軌道に不明瞭な近日点運動が残され、その量は上記の +43 秒分の弧度と大きく異なるものではありませんでした。実験結果の不確実性は数秒程度です。

（b）重力場による光線の偏角

第22節において既に述べたように、一般相対性理論によれば、光線が重力場を通過する際にその経路が曲がるということが予測されます。この曲がり方は、重力場を通って投げ出された物体の経路が経験するものと似ています。この理論に基づいては、太陽から中心までの距離 Δ 太陽半径の位置で通過する光線はその方向に偏角（α）を生じるはずです。

理論によれば、この偏角の半分は太陽のニュートン力学的な引力によって生じ、もう半分は太陽により空間が幾何学的に修正（「曲率」）されることによって生じます。

この結果は、日食中の星の写真撮影によって実験的に検証することができます。唯一の理由として、他の時間帯では太陽からの光により大気が強く照らされ、太陽の近くにある星が見えないためです。予想される効果は添付図から明らかになります。太陽（S）がない場合、地球から観測すると遠くに位置する星は方向D1に見られるべきです。しかし、光の偏角により生じた影響によって、その星は方向D2に見え、すなわち太陽の中心よりやや離れた位置に見えることになります。

実際には、次のように検証されます。日食中に太陽周辺の星を撮影し、その後数ヶ月前後で太陽が他の位置にあるときに同じ星を撮影します。日食写真と標準写真を比較すると、星の位置は太陽中心から放射状に外側へ（太陽の中心方向から離れて）移動し、その量は角度αに対応するはずです。

この重要な推論についての調査には、英国王立学会と王立天文学会に感謝します。第一次世界大戦や戦争によって引き起こされた物的・心理的な困難にも屈することなく、これらの団体はソブラール（ブラジル）と西アフリカのプリンシペ島へ2つの探検隊を送り、1919年5月29日の日食写真を得るためにエディントン、コットンゲム、クロメリン、ダビッドソンといった英国最高峰の天文学者を派遣しました。日食写真と比較写真間の予想される相対的な差異は僅か数ミリメートル程度でしたから、写真撮影およびその後の測定に極めて高い精度が必要でした。

測定結果は理論を完全に満足させる形で確認されました。観測と計算された恒星の偏差（秒弧単位）の直角成分は次の結果表に示されています：



\begin{center}
\begin{tabular}{lrr}
\toprule
星の番号 & 観測された偏角（秒） & 計算された偏角（秒） \\
\midrule
Star 5  & $-0.19$ & $-0.22$ \\
Star 4  & $+0.29$ & $+0.31$ \\
Star 3  & $+0.11$ & $+0.10$ \\
Star 6  & $+0.20$ & $+0.12$ \\
Star 10 & $+0.10$ & $+0.04$ \\
Star 2  & $-0.08$ & $+0.09$ \\
\bottomrule
\end{tabular}
\end{center}




(c) 赤色方向へのスペクトル線の移動

第XXIII節において、回転系K′がGalilei系Kに対して回転している場合、同一構造の時計であっても、回転参照体に対する静止とされる位置にあるときの時計の進歩率は異なることが示されました。現在、この依存関係を定量的に検討します。半径rから中心点まで距離がある時計は、Kに対して相対的な速度vが与えられます：

$v$ = ω$r$,


ここでωは回転盤K′の角速度を表します。時計が静止しているとき、Kに対する単位時間あたりの Tick数（時計の「率」）をv0で表すと、Kに対して相対的に移動しているときの時計の「率」$v$は第XII節に従って与えられます：



\[ \nu = \nu_0 \sqrt{1 - \frac{v^2}{c^2}} \]




または十分な精度では、



\[ \nu = \nu_0 \left(1 - \frac{v^2}{2c^2}\right) \]




この式は以下の形でも表現できます：



\[ \nu = \nu_0 \left(1 + \frac{\phi}{c^2}\right) \]




ここで、回転盤上の位置と中心点間の centrifugal 力による位相差 φ を表すと、つまり、時計のある位置から中心点まで回転盤上を移動するために単位質量に対して反対に働く centrifugal 力に対する作業量（負の値）は、



\[ \phi = -\frac{v^2}{2} = -\frac{r^2 \omega^2}{2} \]




この式から、



\[ \frac{\nu_1 - \nu_2}{\nu_0} = \frac{\phi_1 - \phi_2}{c^2} \]




まず、この式から明らかになるのは、時計が盤の中心から異なる距離にある場合、同一構造の時計でも異なる進歩率で動作することです。これは回転する盤と共に回転している観測者にとっても有効な結果です。

次に、盤上から見れば、後者は位相 φ の重力場に存在していると考えることができます。したがって、得られた結果は一般的にはすべての重力場で成り立ちます。さらに、発光原子を時計とみなすことができるので、以下の声明が成立します：

原子は、その位置にある重力場の位相に依存して吸収または放出する光の周波数を吸収または放出します。

天体表面に存在する原子の周波数は、同じ元素を持つ原子が自由空間（または小さな天体の表面）に存在する場合よりも若干低くなります。

ここで φ = – $K (M/r)$ であり、$K$ はニュートンの重力定数、$M$ は天体の質量です。したがって、星の表面で生成されたスペクトル線と地球の表面で生成された同じ元素のスペクトル線に比べて赤方向への位相移動が発生するはずです。この位相移動の量は



\[ \frac{\Delta \nu}{\nu_0} = -\frac{G M}{c^2 R} \]



太陽の場合、理論によって予測される赤方向への位相移動は約波長の2千万分の一です。恒星の場合、信頼性のある計算は不可能であり、一般に天体の質量 $M$ または半径 $r$ のいずれも知られていないためです。

この効果が存在するかどうかは未解決の問題であり、1920年現在まで天文学者たちはその解明に向けて熱心に取り組んでいます。太陽の場合の効果の小ささから、その存在を判断するのは難しいです。グレーベとバッケム（ボン）、エバシェッドおよびシュワルツシル德による測定結果により、効果の存在はほぼ疑いようがありませんが、他の研究者たちは特にスタインジョンによって測定された結果から逆の結論に至っています。

固定星の統計的調査では、スペクトルの可視光端方向への位相移動が明らかになっていますが、現在利用可能なデータの検討は、これらの移動が実際には重力効果によるものであるかどうかについて明確な結論を導き出すことができません。観測結果は、ここに関心を持った問題から立場に立って詳細に議論されたエドゥアルト・フランキッチの論文「一般相対性理論の検討」（『自然科学』1919年35号、 Julius Springer, Berlin）で集められました。

あらゆる事件には、数年以内に明確な結論が得られるでしょう。スペクトル線の赤方位移が重力ポテンシャルによるものでないならば、一般相対性理論は成り立たなくなります。一方、スペクトル線の位相が重力ポテンシャルによるものであることが明確に判明すれば、その位相を研究することによって天体の質量に関する重要な情報が得られるでしょう。


\clearpage
\chapter{付録 IV：The Structure of Space According to the General Theory of Relativity (Supplementary)}

Appendix IV

宇宙の構造に関する一般相対性理論による考察（第三十二節補遺）

この短い本の最初の版が出版されて以来、大規模な宇宙の構造についての知識（「宇宙論的問題」）は重要な発展を遂げており、その内容は人間向けの簡潔な説明においても触れられるべきものである。

私自身の当初の考察は以下の二つの仮定に基づいていた：

(1) 全ての空間にわたる平均的な物質密度が存在し、どこでも同じであり且つゼロとは異なる。

(2) 空間の大きさ（「半径」）は時間によらず独立している。

これらの仮定は一般相対性理論によって整合性を持つことが示され、しかしそれは引力方程式に仮定的な項を追加することでしか可能であり、その項は理論自体には必要ではなく、また理論的見地から自然とは思われていなかった（「引力方程式の宇宙論的項」）。

当時の私にとって仮定 (2) が避けられないものと感じられていた。なぜなら、それを離れるならば底知れない推測に陥るだろうと考えたからである。

しかし、二十代頃には、ロシアの数学者フリードマンは純粋な理論的見地からは自然な別の仮定を示した。彼は、仮定 (1) を保持しつつ、引力方程式に不自然な宇宙論的項を入れることなく「世界半径」が時間に依存する（空間の拡大）解があることを理解した。つまり、フリードマンによれば、理論は空間の拡大を要求しているという意味で、そのような解があると言える。

数年後には、ハブラーによって特別な調査により、超銀河系（「 Milky Ways 」）から放出されるスペクトル線が距離に比例して赤色シフトを示すことが確認された。このことは現在の知識に基づいてドッペルフェルト原理において解釈されるとすれば、恒星系全体の拡大運動を意味し、これはフリードマンによって引力方程式によって要求されるものである。したがって、ハブラーの発見は理論に対する部分的な確認とみなすことができる。

しかし、奇妙な困難が生じます。ハブラーによって発見された銀河線移動を拡大解釈（理論的な観点から疑いようがない）として解釈すると、この拡大の起源は「ただ」約10億年前にまでさかのぼるとなりますが、物理学的天文学者からは個々の星や星団の発展が著しく長い期間を要することが推測されます。このような不一致をどのように解決するかは全く不明です。

さらに付け加えます。空間の拡大理論と天文学の実証データを合わせると、（三次元）空間の有限性や無限性についての決定は達成されませんが、原始的な「静止」仮説では空間の閉じた（有限な）性質が導かれていました。

$K$ = 座標系

$x, y$ = 二次元座標

$x, y, z$ = 三次元座標

$x, y, z, t$ = 四次元座標

$t$ = 時間

$I$ = 距离

$v$ = 速度

$F$ = 力

$G$ = 重力場


\clearpage

\backmatter

\chapter*{解題：相対性理論の誕生と20世紀物理学の革命}
\addcontentsline{toc}{chapter}{解題：相対性理論の誕生と20世紀物理学の革命}

アインシュタインの思考の軌跡を振り返るとき、私たちが目撃するのは「常識に対する冷徹な批判」と「対称性と不変性への絶対的な確信」です。

特殊相対性理論（1905年）は、光速度不変の原理とガリレイの相対性原理を両立させるために、人類が数千年にわたって自明視してきた「絶対時間」と「絶対空間」の概念を根底から解体しました。動く時計の遅れ、ローレンツ収縮、そして質量とエネルギーの等価性 $E = mc^2$ は、すべてこの幾何学的時空（ミンコフスキー空間）の必然的帰結として演繹されます。

さらに一般相対性理論（1915年）では、「等価原理」を手掛かりとして、ニュートン力学における「遠隔作用としての万有引力」を「物質とエネルギーによって歪められた時空の曲率（リーマン幾何学）」として幾何学的に完全に書き直しました。1919年のエディントンによる皆既日食観測での光線湾曲の確認は、世界中の科学界に衝撃を与え、アインシュタインの名を不滅のものとしました。

本全訳版が、原典の精緻な論理と物理的直観をそのままに味わうための不朽の道標となることを切に願います。

\end{document}
